% Revised November 5, 2026.
\documentclass[11pt]{amsart}

\usepackage[margin=1.15in]{geometry}
\usepackage[T1]{fontenc}
\usepackage{lmodern}
\usepackage{amssymb,mathtools}
\usepackage{microtype}
\usepackage{enumitem,etoolbox}
\usepackage{booktabs,tabularx,placeins}
\usepackage{xcolor}
\usepackage[colorlinks=true,linkcolor=blue!55!black,citecolor=green!35!black,urlcolor=blue!55!black]{hyperref}

\newtheorem{theorem}{Theorem}[section]
\newtheorem{proposition}[theorem]{Proposition}
\newtheorem{lemma}[theorem]{Lemma}
\newtheorem{corollary}[theorem]{Corollary}
\theoremstyle{definition}
\newtheorem{definition}[theorem]{Definition}
\theoremstyle{remark}
\newtheorem{remark}[theorem]{Remark}
\numberwithin{equation}{section}

\newcommand{\C}{\mathbb C}
\newcommand{\R}{\mathbb R}
\newcommand{\PP}{\mathbf P}
\newcommand{\OO}{\mathcal O}
\newcommand{\Tot}{\operatorname{Tot}}
\newcommand{\Bl}{\operatorname{Bl}}
\newcommand{\ddc}{dd^c}
\newcommand{\psh}{\mathrm{psh}}
\newcommand{\Levi}{\mathcal L}
\newcommand{\barp}{\bar\partial}
\newcommand{\supp}{\operatorname{supp}}
\newcommand{\base}{\operatorname{pr}}
\newcommand{\Uone}{\mathrm U(1)}

\title[Regularity of plurisubharmonic exhaustions]{Regularity of
plurisubharmonic exhaustions and related problems}
\author{Zhijian Chen}
\address{Institute of Mathematics, Academy of Mathematics and Systems Science, Chinese Academy of Sciences, Beijing 100190, China}
\email{chenzhijian@amss.ac.cn}
\author{Bo Xiao}
\address{School of Mathematical Sciences, University of Chinese Academy of Sciences, Beijing, China}
\email{xb1994@amss.ac.cn}
\date{November 5, 2026}
\subjclass[2020]{Primary 32F17; Secondary 32Q15, 32L05, 32U40}
\keywords{Plurisubharmonic exhaustion, K\"ahler manifold, semipositive line bundle,
ruled surface, Hironaka modification, Ueda theory, class VII surface, boundary divisor, positive current}

\setcounter{tocdepth}{1}
% Retain normal section spacing after the contents, without amsart's extra skip.
\makeatletter
\patchcmd{\@starttoc}{\addvspace{32\p@\@plus14\p@}}{}{}%
  {\PackageError{psh-exhaustion}{Cannot adjust contents spacing}{}}
\makeatother
\begin{document}

\begin{abstract}
We construct quasi-projective and non-K\"ahler threefolds with locally
Lipschitz plurisubharmonic (psh) exhaustions but no $C^1$ psh exhaustion,
answering Nakano's smoothing question negatively in dimensions at least
three. We prove an affirmative result for divisor complements in compact
K\"ahler surfaces. For line bundles over compact complex manifolds, we
characterize the existence of psh exhaustions of the dual total space by
semipositive metrics of the same $C^r$, H\"older, or real analytic
regularity, establishing the seminegative converse discussed by Ohsawa.
We also prove a holomorphic convexity criterion for reduced complex
spaces fibered by curves over affine space, extending
V\^{\i}j\^{\i}itu's curve-fibering argument to singular spaces and giving
a higher-dimensional analogue of Ohsawa's surface theorem.
\end{abstract}
\maketitle
\tableofcontents

\section{Introduction to the main results and applications}\label{sec:intro}

A \emph{continuous plurisubharmonic (psh) exhaustion} of a reduced complex space
$X$ is a continuous function $\rho:X\to\R$ whose pullback along every
holomorphic disc is subharmonic and whose sublevel sets $\{\rho\le c\}$
are compact for all $c\in\R$. We call $X$ \emph{weakly $1$-complete},
or \emph{weakly pseudoconvex}, if it admits a smooth psh exhaustion.

Can smoothness be replaced by continuity in this definition?
Nakano asked this in 1980, calling manifolds with a continuous
psh exhaustion \emph{Lelong-pseudoconvex}
\cite[\S I.1, Problem~1, p.~162]{Wakabayashi1980}. Demailly's thesis
allows any psh exhaustion in the definition of weak pseudoconvexity
\cite[Introduction, \S1, p.~3]{Demailly1982}. On a K\"ahler manifold,
a continuous psh exhaustion already gives a complete K\"ahler metric,
as needed for his $L^2$ estimates for solving $\barp u=f$
\cite[Theorems~1.3 and~4.1]{DemaillyL2}.

\subsection{Main Results}
In the following theorem we construct a quasi-projective counterexample.

\begin{theorem}\label{thm:main}
There exists a connected smooth quasi-projective K\"ahler threefold
$M$ admitting a psh exhaustion of class $C^{0,1}_{\mathrm{loc}}$
but no psh exhaustion of class $C^{1,0}$.
\end{theorem}

Here by $C^{k,\alpha}_{\mathrm{loc}}$, for $k\ge0$ an integer and
$0<\alpha\le1$, we mean $C^k$ functions whose derivatives of order $k$
are locally $\alpha$-H\"older in real coordinates. Thus
$C^{0,1}_{\mathrm{loc}}$ means locally Lipschitz. At exponent zero
we write $C^{k,0}=C^k$, so $C^{1,0}$ means continuously differentiable.

The construction uses Grauert's ruled-surface example
\cite[\S4, 8(d)]{Grauert1962} and Koike's variant
\cite[Example~4.1]{Koike2021}. In each case the ruled surface $S$ has a
negative section $Y$ and a distinguished line bundle $L$. This bundle
is nef and big (Proposition~\ref{prop:ruled-geometry}) and admits a
locally Lipschitz semipositive metric (Proposition~\ref{prop:continuous-metric}).
By Proposition~\ref{prop:no-metric}, no neighborhood of $Y$ carries
a $C^1$ semipositive metric on $L$. A Hermitian metric $h$ is
semipositive if its weights
$\varphi=-\log\|e\|_h^2$ in local holomorphic frames $e$ are psh.
We take $M=\Tot(L^{-1})$. By the regularity criterion
(Theorem~\ref{cor:criterion}),
\[
\begin{aligned}
&M\text{ has a psh exhaustion of class }\mathcal R
\\&\qquad\Longleftrightarrow\quad
L\text{ has a semipositive metric of class }\mathcal R.
\end{aligned}
\]
This equivalence holds for every holomorphic line bundle over a compact
complex manifold and every $C^r$, local H\"older, or real analytic
($C^\omega$) class $\mathcal R$ in \eqref{eq:regularity-classes}.
The metric obstruction on $L$ therefore rules out every $C^1$ psh
exhaustion of $M$.

Taking the product of $M$ with a Hopf surface $H$ gives a
five-dimensional non-K\"ahler counterexample. Since $H$ is compact,
projection to $M$ pulls its locally Lipschitz psh exhaustion back to
an exhaustion of $M\times H$. A $C^1$ psh exhaustion would restrict
to one on the closed slice $M\times\{h\}$, and a K\"ahler form would
restrict to one on $\{m\}\times H$. To obtain a non-K\"ahler
counterexample in dimension three, we use a modification with a
compact center.

\begin{theorem}\label{thm:nonkahler-threefold}
Let $M$ be the threefold in Theorem~\ref{thm:main}.
For every $x\in M$ there is a smooth non-K\"ahler threefold $N$ and a
proper holomorphic modification $\mu:N\to M$ such that
\[
N\setminus\mu^{-1}(x)\simeq M\setminus\{x\}.
\]
The threefold $N$ admits a psh exhaustion of class $C^{0,1}_{\mathrm{loc}}$
but no psh exhaustion of class $C^{1,0}$.
\end{theorem}

We first blow up $x$ and choose a line and a conic in the exceptional
plane. Following Hironaka's construction \cite[\S6]{Barlet2015}, we
blow up these curves in opposite orders near their two intersection
points. Two exceptional rational curves then have homology classes
whose sum is zero, which rules out a K\"ahler form on $N$.
The locally Lipschitz psh exhaustion of $M$ pulls back to $N$.
Conversely, a $C^1$ psh exhaustion of $N$ could be made constant
near $\mu^{-1}(x)$ by composing it with a smooth convex nondecreasing
function tending to infinity. It would then descend to a $C^1$ psh
exhaustion of $M$ (Proposition~\ref{prop:compact-modification}).

It remains open whether every Lelong-pseudoconvex complex manifold of
dimension two is weakly $1$-complete. The next theorem rules
out counterexamples among complements of divisors in compact
K\"ahler surfaces and, in particular, among smooth quasi-projective
surfaces.

\begin{theorem}\label{thm:surface-introduction}
Let $S$ be a smooth compact connected K\"ahler surface and $D\ne0$
an effective divisor. If $S\setminus\supp D$ admits a continuous
psh exhaustion, it admits a smooth one. Consequently, the same
implication holds for every smooth complex quasi-projective surface.
\end{theorem}

Our proof strategy is to construct a smooth exhaustion of
$X=S\setminus\supp D$ in two cases. If $X$ is a modification of a
Stein space, we pull back a smooth strictly psh exhaustion from that
space. Otherwise we average positive $\partial\barp$-closed currents
supported on levels of an exhaustion. The resulting current has zero
Lelong numbers and compact support in $X$. Using the current theorem
of Dinh--Nguy\^en--Sibony \cite{DNS22} and the Hodge index theorem,
we show that there exists a divisor $A=\sum_i a_iD_i$ with positive
integer coefficients such that
\[
\supp A=\supp D,\qquad A^2=0,\qquad
\OO_S(A)\text{ is unitary flat near }D.
\]
Let $s_A$ be the canonical section of $\OO_S(A)$ and $h$ a flat
Hermitian metric on this bundle near $\supp D$. The function $-\log\|s_A\|_h^2$
is pluriharmonic on a punctured neighborhood of $\supp D$ and tends
uniformly to $+\infty$ on approach to $\supp D$. This is the
\emph{pluriharmonic logarithmic collar} from which we construct the
smooth exhaustion.
Section~\ref{sec:surfaces} also gives separate numerical and
neighborhood criteria for existence of an exhaustion.

For a single smooth connected boundary curve $C\subset S$, the
self-intersection $C^2=\deg N_{C/S}$ and the holomorphic functions
on $X=S\setminus C$ are related by
\[
 C^2<0\quad\Longrightarrow\quad\OO(X)=\C
      \quad\Longrightarrow\quad C^2\le0.
\]
Indeed, when $C^2<0$, Grauert's contraction criterion
\cite[\S4, 8(e)]{Grauert1962} identifies $X$ with the complement of
a point in a compact normal surface, and every holomorphic function
extends across that point. When $C^2>0$, the curve $C$ is nef and
$(K_S-mC)\cdot C<0$ for $m$ large. Serre duality and Riemann--Roch
then give $h^0(S,\OO_S(mC))>1$. Dividing a section independent
of $s_C^m$ by $s_C^m$ gives a nonconstant holomorphic function on $X$.

Ohsawa proved that a connected noncompact complex surface admitting
a smooth psh exhaustion is holomorphically convex precisely when it has
a nonconstant holomorphic function \cite[Proposition~1.4, p.~156]{Ohsawa1981}.
The following analogue for spaces fibered by curves assumes only a
continuous psh exhaustion. For a complement $S\setminus C$ admitting
such an exhaustion, it gives holomorphic convexity when $C^2>0$, and
also when $C^2=0$ if the complement carries a nonconstant
holomorphic function.

\begin{theorem}\label{thm:curve-fiber-introduction}
Let $X$ be a reduced complex space of pure dimension $n\ge2$ with a
continuous real-valued psh exhaustion. If a holomorphic map
\[
F:X\longrightarrow\C^{n-1}
\]
has every nonempty fiber of dimension at most one, then $X$ is
holomorphically convex and admits a smooth psh exhaustion.
\end{theorem}

Analytic fiber cycles and Bishop compactness establish the stability
needed for the convexity argument.

Throughout this paper, \emph{regularization} means finding a smooth
psh exhaustion on a space that admits a continuous one. Among smooth
quasi-projective manifolds, Theorems~\ref{thm:main}
and~\ref{thm:surface-introduction} place the least possible
counterexample dimension at three.

\subsection{Applications to earlier problems}\label{subsec:applications}

\begin{enumerate}[label=\textbf{(\roman*)},widest=iii,leftmargin=*]
\item \textbf{Nakano's smoothing question.}
As recalled above, Nakano asked whether every Lelong-pseudoconvex
manifold is weakly $1$-complete.
Choose a locally Lipschitz psh exhaustion $\rho\ge0$ on the threefold
$M$ of Theorem~\ref{thm:main}.
For each $n>3$, the product $M\times\C^{n-3}$ is quasi-projective and
admits the locally Lipschitz psh exhaustion
\[
 \rho_n(x,z)=\rho(x)+\|z\|^2.
\]
A $C^1$ psh exhaustion on this product would restrict to one on the
closed slice $M\times\{0\}$, contradicting Theorem~\ref{thm:main}.
Thus the answer is negative in every dimension at least three.
Theorem~\ref{thm:surface-introduction} gives an affirmative answer for
surfaces with a divisorial compactification in a compact K\"ahler
surface, including smooth quasi-projective surfaces.

\item \textbf{Filip--Tosatti and the bounded-potential question.}
Filip and Tosatti asked to prove that the distinguished nef and big
line bundle in Grauert's ruled-surface example is not smoothly
semipositive \cite[Problem~2.2, p.~2989]{FilipTosatti2018}.
Zhang proved this for the original example using Koike's first
obstruction class \cite[Theorem~1.2]{Zhang2025}.
He also recorded a question suggested by Tosatti: for the original
or modified Grauert example, does $c_1(L)$ contain a closed positive
current with bounded local potentials and minimal singularities
\cite[\S4]{Zhang2025}?
Proposition~\ref{prop:continuous-metric} and Remark~\ref{rem:minimal}
answer this affirmatively for both examples. The metric has locally
Lipschitz weights and is smooth with strictly positive curvature on
$S\setminus Y$. Its curvature has bounded local potentials and minimal
singularities. In the notation of \eqref{eq:main-line}, the original
and modified examples have
\[
 (Y^2,L^2)=(-1,1),\qquad (Y^2,L^2)=(-2,2),
\]
respectively.

\item \textbf{Koike's first obstruction at \boldmath$C^1$ regularity.}
Koike extended Ueda's obstruction theory to line bundles. His
Theorem~1.4 excludes smooth semipositivity when the first obstruction
class is nonzero, and Example~4.1 applies this to the modified Grauert
example \cite{Koike2021}.
For the Grauert triples $(Y,S,L)$ considered here, the nonzero extension class
$\xi$ in \eqref{eq:extension-intro} gives
\[
 u_1(Y,S,L)=-\xi\ne0
 \quad\text{in }H^1(Y,N_{Y/S}^{-1})\simeq H^1(B,F).
\]
Proposition~\ref{prop:no-metric} shows that no neighborhood of $Y$
carries a $C^1$ Hermitian metric on $L$ with psh weights; the same
holds for every positive tensor power of $L$.
Indeed, Lemma~\ref{lem:first-normal} makes the first normal derivatives
of hypothetical $C^1$ weights holomorphic, and their transition
identities would make $\xi$ a coboundary.

\item \textbf{Ohsawa--V\^{\i}j\^{\i}itu convexity.}
For smooth sources, the proof of Theorem~\ref{thm:curve-fiber-introduction}
follows V\^{\i}j\^{\i}itu's work
\cite[Corollary~4.2 and pp.~463--464]{VijituForum2022}.
Our proof extends the fiber-stability argument to singular total spaces.
For irreducible surfaces, the continuous-exhaustion case already
follows from the stronger theorem
\cite[Theorem~1.1]{Vajaitu2022}.

\item \textbf{Ohsawa's seminegative converse.}
Ohsawa observed that the total space of a seminegative holomorphic
line bundle over a compact complex manifold is weakly $1$-complete
and discussed the converse
\cite[\S6, Proposition~6.1 and the following paragraph]{Ohsawa2021}.
Ohsawa's Theorem~6.1 treats affine line bundles over compact Riemann
surfaces. For linear line bundles, taking $\mathcal R=C^\infty$ in
Theorem~\ref{cor:criterion} proves the converse over any compact complex
manifold. The criterion also preserves the $C^r$, local H\"older, and
real analytic classes in \eqref{eq:regularity-classes}.
\end{enumerate}

\section{Distinguished line bundles on Grauert's ruled surface}\label{sec:metrics}

The two regularity properties of the distinguished bundle $L$ come
from its behavior along the negative section $Y$. A $C^1$ semipositive
metric would split the extension defining the ruled surface. On the
other hand, the negativity of $Y$ allows us to construct a semipositive
metric whose weights are locally Lipschitz across $Y$.

\subsection{Grauert's construction and its line bundle}\label{sec:ruled}

Let $B$ be a smooth connected projective curve of genus $g\ge2$, let
$F$ be a line bundle of degree $d>0$, and choose
$0\ne\xi\in H^1(B,F)$. Let
\begin{equation}\label{eq:extension-intro}
0\longrightarrow F\longrightarrow E\longrightarrow\OO_B
\longrightarrow0
\end{equation}
be an extension with class $\xi$. We write $\PP(E)=\PP_{\mathrm{lines}}(E)$ for the
bundle of lines in $E$, with tautological line bundle $\OO_{\PP(E)}(-1)$.
Set
\begin{equation}\label{eq:main-line}
 S=\PP(E),\qquad \pi:S\to B,\qquad
 Y=\PP(F),\qquad L=\OO_S(Y)\otimes\pi^*F.
\end{equation}
We identify the section $Y$ with $B$ via $\pi$. By GAGA, $E$ is
algebraic and $S$ is a smooth projective surface. The choice $F=K_B$
works for every $g\ge2$, since Serre duality gives
\[
 \deg K_B=2g-2>0,\qquad
 H^1(B,K_B)\simeq H^0(B,\OO_B)^*\simeq\C.
\]
In genus two we may also take $F=\OO_B(p)$ for $p\in B$.
A nonconstant meromorphic function with at most a simple pole at $p$
would give a degree-one map $B\to\PP^1$. Thus $h^0(B,\OO_B(p))=1$,
and Riemann--Roch gives $h^1(B,\OO_B(p))=1$ as well. These two
degree choices recover Grauert's example for $F=\OO_B(p)$
\cite[\S4, 8(d)]{Grauert1962} and Koike's variant for $F=K_B$
\cite[Example~4.1]{Koike2021}.

Choose local frames $f_i$ of $F$ on $U_i\subset B$ and lifts
$v_i$ of $1\in\OO_B$, with
\[
 f_j=g_{ij}f_i,\qquad v_j=v_i+c_{ij}f_i.
\]
The cocycle $\{v_j-v_i=c_{ij}f_i\}$ represents $\xi$.
On $S\setminus Y$, write a line as $[t_if_i+v_i]$. Near $Y$,
write it as $[f_i+w_iv_i]$, with $w_i=1/t_i$ where both coordinates
are defined. Comparing
coefficients on overlaps gives
\begin{equation}\label{eq:affine-transition}
 t_i=g_{ij}t_j+c_{ij},\qquad
 w_i=\frac{w_j}{g_{ij}+c_{ij}w_j},\qquad
 w_j=\frac{g_{ij}w_i}{1-c_{ij}w_i}.
\end{equation}
The quotient map $E\to\OO_B$ induces a section of
$\OO_S(1)$ whose zero divisor is $Y$. Consequently
\begin{equation}\label{eq:restriction}
 \begin{gathered}
 \OO_S(Y)\simeq\OO_S(1),\qquad
 N_{Y/S}\simeq\operatorname{Hom}(F,E/F)\simeq F^{-1},\\
 Y^2=-d,\qquad L|_Y\simeq F^{-1}\otimes F\simeq\OO_Y.
 \end{gathered}
\end{equation}
Let $s_Y$ be the canonical section of $\OO_S(Y)$. Near $Y$,
let $d_i$ be the frame defined by $s_Y=w_id_i$,
and put $\ell_i=d_i\otimes\pi^*f_i$. Then
\begin{equation}\label{eq:L-transition}
 d_j=\frac{1-c_{ij}w_i}{g_{ij}}d_i,\qquad
 \ell_j=(1-c_{ij}w_i)\ell_i.
\end{equation}
In particular, the frames $\ell_i|_Y$ agree and give the
trivialization in \eqref{eq:restriction}.

\begin{proposition}\label{prop:ruled-geometry}
The line bundle $L$ is nef and big, with
\[
 L^2=d,\qquad L\cdot Y=0,\qquad
 \operatorname{Null}(L)=Y.
\]
The bundle $A=L^{\otimes2}\otimes\OO_S(-Y)$ is ample.
Every section of a positive tensor power of $L$ vanishes along $Y$.
In particular, $L$ is not semiample.
\end{proposition}
\begin{proof}
Let $P$ be a fiber of $\pi$. The intersection identities are
\[
 Y^2=-d,\quad Y\cdot\pi^*F=d,\quad (\pi^*F)^2=0,
 \quad Y\cdot P=1,\quad \pi^*F\cdot P=0.
\]
Thus $L^2=d$, $L\cdot Y=0$, and $L\cdot P=1$.
For an irreducible horizontal curve $D\ne Y$, let
$n=\deg(\pi|_{\widetilde D})>0$, where $\widetilde D$ is its
normalization. Then
\[
 L\cdot D=Y\cdot D+nd>0.
\]
Thus $L$ is nef, and $L^2=d>0$ makes it big by the numerical
criterion for nef bundles on a projective surface. The only curve on
which $L$ has degree zero is $Y$. For $A$ we have
\[
 \begin{gathered}
 A=\OO_S(Y)\otimes\pi^*(F^{\otimes2}),\qquad A^2=3d,\\
 A\cdot Y=d,\qquad A\cdot P=1,\qquad
 A\cdot D=Y\cdot D+2nd>0.
 \end{gathered}
\]
The Nakai--Moishezon criterion now gives ampleness.

To prove the assertion about sections, we show that $S\setminus Y$
contains no compact irreducible curve. Such a curve $D$ would have
to dominate $B$, since each fiber in the complement is an affine
line. Its normalization would give a finite map $p:\widetilde D\to B$
of degree $n>0$. Away from the branch values,
put
\[
 a_i(b)=\frac1n\sum_{x\in p^{-1}(b)}t_i(x).
\]
The sum is holomorphic and single-valued, and is locally bounded
at branch values by compactness of the inverse image of a small
closed disc. It extends holomorphically across them. Equation
\eqref{eq:affine-transition} implies
\[
 a_i=g_{ij}a_j+c_{ij},\qquad
 c_{ij}f_i=a_if_i-a_jf_j.
\]
This makes the cocycle representing $\xi$ a coboundary, a
contradiction.

Now let $\sigma\in H^0(S,L^{\otimes m})$, $m>0$.
Its restriction to $Y$ is constant. If that constant were nonzero,
the zero divisor of $\sigma$ would be disjoint from $Y$ and would
be nonempty, since
\[
 L^{\otimes m}|_P\simeq\OO_{\PP^1}(m),\qquad
 \deg\operatorname{div}(\sigma|_P)=m>0.
\]
It would be a compact curve in $S\setminus Y$, which is
impossible. Hence $\sigma|_Y=0$ for every $m>0$.
\end{proof}

\subsection{The first normal derivative and the $C^1$ obstruction}

On $Y$, the weights of a semipositive metric on $L$ glue to a
subharmonic function and hence are constant in the frames $\ell_i$.
The next lemma determines what plurisubharmonicity then imposes on
their first derivatives normal to $Y$.

\begin{lemma}[First normal derivative]\label{lem:first-normal}
Let $D\subset\C$ be a disc and let $\varphi$ be a real-valued
$C^1$ psh function on a neighborhood of $D\times\{0\}$ in
$\C^2$. If $\varphi(z,0)$ is constant, then
$b(z)=\partial_w\varphi(z,0)$ is holomorphic.
\end{lemma}
\begin{proof}
Subtract the constant $\varphi(z,0)$. By continuity of
$\partial_w\varphi$, the first-order expansion
\begin{equation}\label{eq:first-normal-expansion}
 \varphi(z,w)=2\operatorname{Re}(b(z)w)+o(|w|)
\end{equation}
is uniform on each compact subdisc in the $z$ variable.
Fix a holomorphic function $f$ on $D$. On each subdisc $D_0\Subset D$
and for all sufficiently small $t>0$, the graph restrictions
\[
 u_t^\pm(z)=\frac{\varphi(z,\pm tf(z))}{t}
\]
are subharmonic and converge uniformly to
$\pm2\operatorname{Re}(bf)$ by \eqref{eq:first-normal-expansion}.
Both limits are subharmonic, so $\operatorname{Re}(bf)$ is harmonic
on $D$. Taking $f=1,i,z,iz$ shows that $b$ and $zb$ are harmonic.
It follows that
\[
 0=\Delta(zb)=z\Delta b+4\partial_{\bar z}b.
\]
Thus $\partial_{\bar z}b=0$.
\end{proof}

\begin{proposition}\label{prop:no-metric}
No neighborhood of $Y$ carries a $C^1$ Hermitian metric on $L$
with psh weights. The same holds for $L^{\otimes m}$, $m\ge1$.
\end{proposition}
\begin{proof}
For a Hermitian metric $h$ on a line bundle $A$, its weights
$\varphi_i=-\log\|e_i\|_h^2$ in local frames $e_j=g_{ij}e_i$
satisfy
\begin{equation}\label{eq:metric-transition}
 \varphi_j=\varphi_i-\log|g_{ij}|^2.
\end{equation}
With $d^c=\tfrac{i}{4\pi}(\barp-\partial)$ and
$\ddc=\tfrac{i}{2\pi}\partial\barp$, the curvature is
$\Theta_h(A)=\ddc\varphi_i$, in the sense of currents for weights
of regularity below $C^2$.
Suppose the weights $\varphi_i$ in the frames $\ell_i$ are $C^1$
and psh near $Y$. Applying \eqref{eq:metric-transition} to
\eqref{eq:L-transition} gives
\begin{equation}\label{eq:first-order-weight}
 \varphi_j(z,w_j)=\varphi_i(z,w_i)-\log|1-c_{ij}w_i|^2.
\end{equation}
The restrictions $\varphi_i|_Y$ glue to a subharmonic function
on the compact connected curve $Y$, hence to a constant. Rescale
the metric so that $\varphi_i(z,0)=0$. By
Lemma~\ref{lem:first-normal}, the functions
\[
 b_i(z)=\partial_{w_i}\varphi_i(z,0)
\]
are holomorphic. Differentiating \eqref{eq:first-order-weight} at
$w_i=0$ and using \eqref{eq:affine-transition} gives
\[
 g_{ij}b_j=b_i+c_{ij},\qquad
 v_j-b_jf_j=v_i+(c_{ij}-g_{ij}b_j)f_i=v_i-b_if_i.
\]
Therefore
\[
 s_i:\OO_B|_{U_i}\longrightarrow E|_{U_i},\qquad
 a\longmapsto a(v_i-b_if_i)
\]
glue to a holomorphic right inverse of $E\to\OO_B$,
contradicting $\xi\ne0$. A $C^1$ semipositive metric on
$L^{\otimes m}$ would give one on $L$ by dividing its weights
by $m$.
\end{proof}

Under the identification $N_{Y/S}^{-1}\simeq F$, the local generator
$dw_i|_Y$ corresponds to $f_i$. Thus the first obstruction class is
\[
 u_1(Y,S,L)=\bigl[\{-c_{ij}\,dw_i\}\bigr]=-\xi
 \quad\text{in }H^1(Y,N_{Y/S}^{-1})\simeq H^1(B,F).
\]
Koike's obstruction theorem applies to smooth weights
\cite[Theorem~1.4]{Koike2021}. Wu also computes the smooth
obstruction in Grauert's example \cite[Example~3.7]{Wu2026}.
Lemma~\ref{lem:first-normal} gives the obstruction at $C^1$ regularity.

\subsection{Cusp correction and a global metric}\label{sec:cusp}

To construct the locally Lipschitz metric, we add a multiple of the
norm of the defining section to a smooth local weight. Away from the
curve this term supplies positive curvature in both the tangential and
normal directions.

\begin{lemma}[Cusp correction]\label{lem:cusp}
Let $D$ be a compact smooth connected curve in a complex surface $X$, with
$D^2<0$, and let $H$ be a line bundle near $D$ whose restriction
admits a smooth flat metric. There are smooth metrics $h_0$ on
$H$ and $k$ on $\OO_X(D)$, a neighborhood $U$ of $D$, and
$\Lambda>0$ such that
\[
 h_{\mathrm{loc}}=h_0e^{-\Lambda s},\qquad
 s=\|s_D\|_k,
\]
has locally Lipschitz psh weights on $U$ and smooth strictly psh
weights on $U\setminus D$.
\end{lemma}
\begin{proof}
Extend the flat metric on $H|_D$ to a smooth metric $h_0$ near $D$
and put $\theta=\Theta_{h_0}(H)$. Then $\theta|_{TD}=0$.
Since $\deg N_{D/X}=D^2<0$, choose a metric on $N_{D/X}$
with strictly negative curvature and extend it to a smooth metric
$k$ on $\OO_X(D)$. Indeed, if $\gamma$ is an initial curvature
and $\omega_D$ is an area form of integral one, the equation
\[
 \ddc v=D^2\omega_D-\gamma,\qquad
 \int_D(D^2\omega_D-\gamma)=0
\]
is solvable on the compact curve. Multiplication by $e^{-v}$
gives the prescribed negative curvature. Smooth extensions follow
by extending the logarithm of the ratio of two metrics. Put
\[
 \eta=-\Theta_k(\OO_X(D)),\qquad
 q=\|s_D\|_k^2,\qquad s=q^{1/2}.
\]
In coordinates $D=\{w=0\}$, with $s_D=wd$ and
$\|d\|_k^2=e^\beta$, one has
\begin{equation}\label{eq:baseline-metric}
 q=|w|^2e^\beta,\qquad s=|w|e^{\beta/2},\qquad
 \eta=\ddc\beta,\qquad \eta|_{TD}>0.
\end{equation}
Shrink finitely many relatively compact charts covering $D$ until
$|w\beta_w|\le1/2$. Define smooth vector fields
\[
 H_1=\partial_z-\frac{w\beta_z}{1+w\beta_w}\partial_w,
 \qquad H_2=\partial_w.
\]
They form a frame, with $H_1|_D=\partial_z$, and on $w\ne0$
\[
 \partial\log q(H_1)=0,\qquad
 \partial\log q(H_2)=\frac1w+\beta_w.
\]
The tensor identity
\begin{equation}\label{eq:scaled-cusp}
 \ddc s=\frac{s}{2}\eta+
 \frac{s}{4}\frac{i}{2\pi}
       \partial\log q\wedge\barp\log q
\end{equation}
therefore gives, in the frame $(H_1,H_2)$,
\[
 \begin{aligned}
 (\ddc s)_{1\bar1}&=\tfrac{s}{2}\eta_{1\bar1},&
 (\ddc s)_{1\bar2}&=\tfrac{s}{2}\eta_{1\bar2},\\
 (\ddc s)_{2\bar2}&=\tfrac{s}{2}\eta_{2\bar2}
           +\tfrac{s}{4}|w^{-1}+\beta_w|^2.
 \end{aligned}
\]
These coefficients are tensor evaluations in a smooth frame,
with the common factor $i/(2\pi)$ omitted.
Choose uniform constants $\mu>0$, $A\ge1$ with
\[
 \eta_{1\bar1}\ge\mu,\qquad
 |\eta_{1\bar2}|+|\eta_{2\bar2}|\le A,\qquad
 A^{-1}\le e^\beta\le A.
\]
Since $|1+w\beta_w|\ge1/2$,
\[
 (\ddc s)_{2\bar2}\ge-\frac A2s+
        \frac{s}{16|w|^2}
 \ge-\frac A2s+\frac1{16As}\ge\frac1{32As}
\]
after imposing $s^2\le1/(16A^2)$. Smoothness and
$\theta_{1\bar1}|_D=0$ also give
$|\theta_{1\bar1}|\le C|w|\le C A^{1/2}s$.
Consequently there are $a,b,K>0$, independent of the chart, such
that
\begin{equation}\label{eq:cusp-estimates}
 \begin{aligned}
 (\ddc s)_{1\bar1}&\ge as,&
 (\ddc s)_{2\bar2}&\ge b/s,&
 |(\ddc s)_{1\bar2}|&\le Ks,\\
 \theta_{1\bar1}&\ge-Ks,&
 \theta_{2\bar2}&\ge-K,&
 |\theta_{1\bar2}|&\le K.
 \end{aligned}
\end{equation}
First fix $\Lambda$ so that
\begin{equation}\label{eq:cusp-coefficient}
 \Lambda a-K>0,\qquad
 (\Lambda a-K)\frac{\Lambda b}{2}>(K+1)^2.
\end{equation}
Then shrink $U$ until $Ks\le\Lambda b/2$ and
$\Lambda Ks\le1$. For $T=\theta+\Lambda\ddc s$,
\[
 T_{1\bar1}\ge(\Lambda a-K)s>0,\qquad
 T_{2\bar2}\ge\frac{\Lambda b}{2s}>0,\qquad
 |T_{1\bar2}|\le K+1,
\]
and hence
\[
 \det T\ge(\Lambda a-K)\frac{\Lambda b}{2}-(K+1)^2>0.
\]
Thus $T$ is strictly positive on $U\setminus D$. If $\varphi_0$ is a
local weight of $h_0$, the weight
$\varphi_0+\Lambda s$ is continuous across $D$ and psh off $D$, so
the removable-singularity theorem makes it psh on the whole chart. Finally,
$s=|w|e^{\beta/2}$ is locally Lipschitz, proving the assertion.
\end{proof}

\begin{proposition}[Global cusp metric]\label{prop:global-cusp}
Let $X$ be a smooth projective surface, $D\subset X$ a smooth
connected curve with $D^2<0$, and $H$ a line bundle such that
$H|_D$ admits a smooth flat metric. Suppose
\[
 A=H^{\otimes m}\otimes\OO_X(-a_0D)
\]
is ample for some integers $m,a_0>0$. Then $H$ admits a
locally Lipschitz semipositive Hermitian metric, smooth with
strictly positive curvature on $X\setminus D$.
\end{proposition}
\begin{proof}
Extend the smooth metrics $h_0$ and $k$ from Lemma~\ref{lem:cusp}
to $X$, and set $q=\|s_D\|_k^2$ and $s=q^{1/2}$ globally.
Write $\theta=\Theta_{h_0}(H)$ and $u_{\mathrm{loc}}=\Lambda s$.
Let $h_A$ be a smooth metric
with $\omega_A=\Theta_{h_A}(A)>0$. The identity of line bundles
$H^{\otimes m}=A\otimes\OO_X(a_0D)$ defines a global smooth
function $G$ by
\[
 h_A\otimes k^{a_0}=h_0^m e^{-G}.
\]
For the canonical section $s_D$ of $\OO_X(D)$, the
Poincar\'e--Lelong formula gives
\begin{equation}\label{eq:PL}
 \ddc\log\|s_D\|_k^2=[D]-\Theta_k(\OO_X(D)).
\end{equation}
Hence the divisor metric $h_{[D]}=ke^{-\log q}$ has curvature $[D]$.
Thus the auxiliary metric and its weight relative to $h_0$ are
\begin{equation}\label{eq:log-curvature}
 \begin{gathered}
 h_*^m=h_A\otimes h_{[D]}^{a_0},\qquad
 h_*=h_0e^{-\psi},\qquad
 \psi=\frac{a_0}{m}\log q+\frac Gm,\\
 \theta+\ddc\psi=\frac1m\omega_A+\frac{a_0}{m}[D].
 \end{gathered}
\end{equation}
In particular, $\psi$ is strictly $\theta$-psh off $D$ and
tends uniformly to $-\infty$ on approach to $D$.

Fix $\varepsilon>0$. Choose an even nonnegative smooth function
$\chi_\varepsilon$ supported in $(-\varepsilon,\varepsilon)$
with integral one, and put
\[
 \rho_\varepsilon(t)=\int_\R|t-a|\chi_\varepsilon(a)\,da,
 \qquad
 \mathcal M_\varepsilon(x,y)=
       \frac{x+y+\rho_\varepsilon(x-y)}2.
\]
Then $\rho_\varepsilon''\ge0$, $|\rho_\varepsilon'|\le1$,
and $\rho_\varepsilon(t)=|t|$ for $|t|\ge\varepsilon$.
Consequently $\mathcal M_\varepsilon$ is convex,
nondecreasing in both variables, and
\[
 \mathcal M_\varepsilon(x+c,y+c)=c+\mathcal M_\varepsilon(x,y),
 \qquad
 \mathcal M_\varepsilon(x,y)=\max\{x,y\}\quad
                  (|x-y|\ge\varepsilon).
\]
Writing $\Levi_f(v)=\sum_{j,k} f_{z_j\bar z_k}v_j\overline{v_k}$
for the Levi form of a $C^2$ function $f$, we obtain for smooth
functions $p_1,p_2$
\[
 \Levi_{\mathcal M_\varepsilon(p_1,p_2)}(v)
 =\sum_{j=1}^2(\mathcal M_\varepsilon)_j\Levi_{p_j}(v)
  +\sum_{j,k=1}^2(\mathcal M_\varepsilon)_{jk}
       \partial p_j(v)\overline{\partial p_k(v)}.
\]
The second term is nonnegative, and the coefficients in the
first term are nonnegative and sum to one. Thus strict
plurisubharmonicity is preserved.

Choose $b>0$ so that $\{q\le2b\}\subset U$, where the local
metric is defined, and then choose $N\in\R$ and $0<a<b/4$
such that
\begin{equation}\label{eq:metric-collars}
 \begin{aligned}
 \psi+N&>u_{\mathrm{loc}}+2\varepsilon
       &&\text{on }\{b/2\le q\le b\},\\
 \psi+N&<u_{\mathrm{loc}}-2\varepsilon
       &&\text{on }\{0<q\le2a\}.
 \end{aligned}
\end{equation}
First take $N$ large on the compact shell to obtain the first
inequality, then take $a$ small enough that
\[
 \frac{a_0}{m}\log(2a)+\max_X(G/m)+N<-2\varepsilon.
\]
Define
\begin{equation}\label{eq:global-metric}
 h=h_0e^{-u},\qquad
 u=\begin{cases}
 u_{\mathrm{loc}},&q\le a,\\
 \mathcal M_\varepsilon(u_{\mathrm{loc}},\psi+N),&a<q<b,\\
 \psi+N,&q\ge b.
 \end{cases}
\end{equation}
The expressions coincide on collars of $q=a$ and $q=b$ by
\eqref{eq:metric-collars}. On the annulus, a local weight is
\[
 \varphi_0+u=\mathcal M_\varepsilon
  (\varphi_0+u_{\mathrm{loc}},\varphi_0+\psi+N),
\]
so it is smooth and strictly psh. The inner weight is psh by
Lemma~\ref{lem:cusp}, and the outer one is strictly psh by
\eqref{eq:log-curvature}. The metric is locally Lipschitz near
$D$ and smooth with strictly positive curvature elsewhere.
\end{proof}

\begin{proposition}\label{prop:continuous-metric}\label{sec:global}
For the ruled surface \eqref{eq:main-line}, $L$ admits a locally
Lipschitz semipositive Hermitian metric, smooth with strictly
positive curvature on $S\setminus Y$. Its curvature has bounded
local potentials and minimal singularities.
\end{proposition}
\begin{proof}
Equation \eqref{eq:restriction} and
Proposition~\ref{prop:ruled-geometry} give
\[
 Y^2=-d<0,\qquad L|_Y\simeq\OO_Y,\qquad
 L^{\otimes2}\otimes\OO_S(-Y)\text{ ample}.
\]
Apply Proposition~\ref{prop:global-cusp} with
$(X,D,H,m,a_0)=(S,Y,L,2,1)$. Write the resulting curvature as
$T=\theta+\ddc u\ge0$, where $u$ is continuous on compact $S$.
For every global quasi-psh function $v\not\equiv-\infty$ with
$\theta+\ddc v\ge0$,
\[
 v\le u+\sup_Sv-\inf_Su.
\]
This is the defining comparison for minimal singularities, and
the local potentials $\varphi_0+u$ are bounded.
\end{proof}

\begin{remark}\label{rem:minimal}
The choices $g=2$, $F=\OO_B(p)$ and $F=K_B$ yield respectively
\[
 (Y^2,L^2)=(-1,1),\qquad (Y^2,L^2)=(-2,2).
\]
Proposition~\ref{prop:continuous-metric} gives bounded potentials
with minimal singularities for both classes considered in
\cite[Section~4]{Zhang2025}. Proposition~\ref{prop:no-metric}
excludes $C^1$ semipositive weights in either case.
\end{remark}

\section{The regularity criterion}\label{sec:radius}

For a semipositive Hermitian metric on $A$, the squared norm on the
dual bundle gives a psh exhaustion. To recover a metric from an
exhaustion, we first average over rotations of the fibers. A suitable
sublevel set then meets each fiber in a disc whose radius determines
the metric. We carry out both constructions for the local regularity
classes
\begin{equation}\label{eq:regularity-classes}
\begin{split}
 \mathcal R\in{}&\{C^r:r\in\mathbb Z_{\ge0}\cup\{\infty\}\}
 \cup\{C^\omega\}\\
 &{}\cup\{C^{k,\alpha}_{\mathrm{loc}}:
          k\in\mathbb Z_{\ge0},\ 0<\alpha\le1\}.
\end{split}
\end{equation}
For every class in \eqref{eq:regularity-classes}, the positive metric
coefficients $h_i$, their reciprocals, and the weights
$-\log h_i$ have the same regularity.

\begin{lemma}\label{lem:norm}
Let $A\to Z$ be a holomorphic line bundle over a compact complex
manifold. If $h$ is a semipositive Hermitian metric of class
$\mathcal R$, then $\rho(v)=\|v\|_{h^{-1}}^2$ is a psh exhaustion
of $\Tot(A^{-1})$ of the same class.
\end{lemma}
\begin{proof}
Let $e_i$ be a local frame of $A$ over $B_i\subset Z$, with
$\|e_i\|_h^2=e^{-\varphi_i}$. For $v=\lambda_i e_i^*$,
\begin{equation}\label{eq:local-norm}
 \rho(z,\lambda_i)=|\lambda_i|^2e^{\varphi_i(z)}.
\end{equation}
This formula gives the claimed regularity on the total space.
Off the zero section, $\log\rho=\log|\lambda_i|^2+\varphi_i$ is psh,
so $\rho$ is psh there by convex increasing composition.
Continuity and removability across the zero section give
plurisubharmonicity everywhere.

Fix a smooth Hermitian metric $k$ on $A^{-1}$. The ratio of $h^{-1}$
to $k$ is a positive continuous function on the compact base. Its
minimum $m>0$ gives $\rho(v)\ge m\|v\|_k^2$, so each sublevel of
$\rho$ is a closed subset of a compact disc bundle.
\end{proof}

\begin{lemma}[Compact-group averaging]\label{lem:average}
If a compact group $G$ acts continuously on a locally compact
Hausdorff space and $F_0$ is a continuous exhaustion, then
$F(x)=\int_GF_0(gx)\,dg$, with normalized Haar measure $dg$, is
an exhaustion. For the fiberwise circle
action on a holomorphic line bundle, averaging preserves
plurisubharmonicity and every class in \eqref{eq:regularity-classes}.
\end{lemma}
\begin{proof}
Continuity follows from uniform continuity on compact orbit sets.
For $K_c=\{F_0\le c\}$,
\[
 \{F\le c\}\subset G\cdot K_c.
\]
Indeed, $x\notin G\cdot K_c$ implies
$\min_{g\in G}F_0(gx)>c$, and therefore $F(x)>c$.
The set $G\cdot K_c$ is the continuous image of the compact set
$G\times K_c$, so its closed subset $\{F\le c\}$ is compact.

For the circle action $R_\theta(v)=e^{i\theta}v$, put
\begin{equation}\label{eq:circle-average}
 F(v)=\frac1{2\pi}\int_0^{2\pi}F_0(R_\theta v)\,d\theta.
\end{equation}
If $F_0$ is psh, integration of the submean inequality on each
holomorphic disc gives the same inequality for $F$. Invariance
under rotation follows from translation invariance of $d\theta$.

In real bundle coordinates $x=(z,\lambda)$, write
$T_\theta x=(z,e^{i\theta}\lambda)$. The maps $T_\theta$ are linear
with uniformly bounded norms. For $F_0\in C^r$, $r\ge1$,
differentiation on compact coordinate sets yields
\[
 D^\beta F(x)=\frac1{2\pi}\int_0^{2\pi}
 D_x^\beta(F_0\circ T_\theta)(x)\,d\theta,
 \qquad |\beta|\le r.
\]
For $F_0\in C^{k,\alpha}_{\mathrm{loc}}$, a finite cover of the
compact orbit set gives, uniformly in $\theta$,
\[
 |D^\beta(F_0\circ T_\theta)(x)
   -D^\beta(F_0\circ T_\theta)(y)|
 \le C|x-y|^\alpha,\qquad |\beta|=k.
\]
Integration gives the same estimate for $D^\beta F$.

For the analytic case, $\mathcal G(x,\theta)=F_0(T_\theta x)$ is
jointly real analytic near the compact set $\{x_0\}\times S^1$.
A finite cover by analytic coordinate neighborhoods and Cauchy's
estimates give constants $M,R>0$ such that
\begin{equation}\label{eq:analytic-average-estimate}
 \left|\frac{D_x^\beta\mathcal G(x_0,\theta)}{\beta!}\right|
 \le MR^{-|\beta|}
 \qquad(\beta\in\mathbb Z_{\ge0}^{m},\ \theta\in S^1),
\end{equation}
where $m$ is the real coordinate dimension. The Taylor series in $x$
converges uniformly in $\theta$ on every smaller polydisc. Integrating
term by term gives
\[
 F(x)=\sum_\beta\left(\frac1{2\pi}\int_0^{2\pi}
      \frac{D_x^\beta\mathcal G(x_0,\theta)}{\beta!}\,d\theta\right)
      (x-x_0)^\beta.
\]
This convergent power series proves $F\in C^\omega$.
\end{proof}

\begin{proposition}[Recovery of the metric]\label{prop:radius}
Let $A\to Z$ be a holomorphic line bundle over a compact complex
manifold. If $\Tot(A^{-1})$ admits a psh exhaustion of class
$\mathcal R$ in \eqref{eq:regularity-classes}, then $A$ admits a
semipositive Hermitian metric of the same class.
\end{proposition}
\begin{proof}
Let $F$ be the circle average \eqref{eq:circle-average} of the given
exhaustion. For a local frame $e_i$ of $A$ over $B_i$, define
\[
 U_i(z,t)=F(e^{t/2}e_i^*(z)),\qquad
 V_i(z,\zeta)=F(e^{\zeta/2}e_i^*(z)).
\]
Both parametrizations are real analytic, so $U_i$ and $V_i$
inherit the class $\mathcal R$. The second parametrization is
holomorphic, so $V_i$ is psh. Circle invariance gives
\[
 V_i(z,\zeta)=U_i(z,\operatorname{Re}\zeta),\qquad
 F(\lambda_i e_i^*(z))=U_i(z,\log|\lambda_i|^2)
 \quad(\lambda_i\ne0).
\]
Fix $z$ and write $u(t)=U_i(z,t)$. Since $V_i$ is independent of
$\operatorname{Im}\zeta$, its subharmonicity in $\zeta$ gives
$u''\ge0$ in the sense of distributions. Thus $u$ is convex.
Continuity at the zero section and
properness of $F$ give
\begin{equation}\label{eq:radial-limits}
 \lim_{t\to-\infty}u(t)=F(0_z),\qquad
 \lim_{t\to+\infty}u(t)=+\infty.
\end{equation}
If $u(b)<u(a)$ for $a<b$, the negative secant slope
$d=(u(b)-u(a))/(b-a)$ would give
\[
 \frac{u(a)-u(x)}{a-x}\le d<0,\qquad
 u(x)\ge u(a)-d(a-x)\longrightarrow+\infty
 \quad(x\to-\infty),
\]
contrary to \eqref{eq:radial-limits}. Thus $u$ is nondecreasing.

Choose $c>\max_{z\in Z}F(0_z)$. The intermediate value theorem
gives a point $s$ with $u(s)=c$. If $u(s)=u(t)=c$ for $s<t$, then
for $x<s$ the ordered secant slopes give
\[
 0\le\frac{u(s)-u(x)}{s-x}
 \le\frac{u(t)-u(s)}{t-s}=0.
\]
This would give $u(x)=c$ for all $x<s$, again contradicting
\eqref{eq:radial-limits}. Consequently there is a unique
$\sigma_i(z)$ such that
\begin{equation}\label{eq:implicit-radius}
 U_i(z,\sigma_i(z))=c,
 \qquad U_i(z,t)\le c\ \Longleftrightarrow\ t\le\sigma_i(z).
\end{equation}
In particular,
\begin{equation}\label{eq:radius-sublevel}
 \{F<c\}\cap A^{-1}|_{B_i}
 =\{(z,\lambda_i):|\lambda_i|^2<e^{\sigma_i(z)}\}.
\end{equation}
For $s_0=\sigma_i(z_0)$ and $\varepsilon>0$,
\[
 U_i(z_0,s_0-\varepsilon)<c<U_i(z_0,s_0+\varepsilon).
\]
These strict inequalities persist near $z_0$ and force
$s_0-\varepsilon<\sigma_i(z)<s_0+\varepsilon$. Thus $\sigma_i$
is continuous.

To prove that $-\sigma_i$ is psh, let $f$ be a holomorphic map from a
neighborhood of $\overline\Delta$ to $B_i$. Let $H$ be the harmonic
extension of $\sigma_i\circ f|_{\partial\Delta}$, continuous on
$\overline\Delta$. Since $\Delta$ is simply
connected, there is a holomorphic $g$ with $\operatorname{Re}g=H$.
The function
\[
 q(\zeta)=F\bigl(e^{g(\zeta)/2}e_i^*(f(\zeta))\bigr)
          =U_i(f(\zeta),H(\zeta))
\]
is subharmonic on $\Delta$, and the second expression extends
continuously to the boundary with value $c$. The maximum principle
and \eqref{eq:implicit-radius} give
\[
 q\le c\quad\Longrightarrow\quad
 H\le\sigma_i\circ f\quad\Longrightarrow\quad
 -\sigma_i(f(0))\le\frac1{2\pi}\int_0^{2\pi}
       -\sigma_i(f(e^{i\theta}))\,d\theta.
\]
The submean inequality for every such disc proves that $-\sigma_i$ is psh.

If $e_j=g_{ij}e_i$, then $\lambda_j=g_{ij}\lambda_i$, and the
radii of the same sublevel in \eqref{eq:radius-sublevel} satisfy
\[
 e^{\sigma_j}=|g_{ij}|^2e^{\sigma_i},\qquad
 \sigma_j=\sigma_i+\log|g_{ij}|^2.
\]
Thus
\begin{equation}\label{eq:radius-metric}
 \psi_i=-\sigma_i,\qquad
 \psi_j=\psi_i-\log|g_{ij}|^2,\qquad
 \|e_i\|_h^2=e^{-\psi_i}=e^{\sigma_i}
\end{equation}
defines a continuous Hermitian metric on $A$ with
$\Theta_h(A)=\ddc\psi_i\ge0$.

We now recover the regularity of $F$. For $C^1$ data,
convexity gives, for any $a<\sigma_i(z)$,
\begin{equation}\label{eq:radius-transversality}
 U_{i,t}(z,\sigma_i(z))
 \ge\frac{c-U_i(z,a)}{\sigma_i(z)-a}>0.
\end{equation}
The implicit function theorem gives $\sigma_i\in C^r$ whenever
$F\in C^r$, for $1\le r\le\infty$. Its real analytic version gives
the same conclusion for $C^\omega$.

For $C^{0,\alpha}_{\mathrm{loc}}$ data, omit the index $i$ and
choose
\[
 a<b<\sigma(z_0)<d,\qquad U(z_0,a)<U(z_0,b).
\]
The limits \eqref{eq:radial-limits} permit this choice.
By continuity, there are a relatively compact base ball $B$
around $z_0$ and $m>0$ such that
\[
 b<\sigma(z)<d,\qquad
 \frac{U(z,b)-U(z,a)}{b-a}\ge m
 \qquad(z\in\overline B).
\]
The local H\"older regularity of $U$ on a neighborhood of
$\overline B\times[b,d]$ gives one constant $C$ such that
\[
 |U(z,t)-U(w,t)|\le C|z-w|^\alpha
 \qquad(z,w\in B,\ t\in[b,d]).
\]
Ordered secant slopes then yield
\[
 U(z,t)-U(z,s)\ge m(t-s),\qquad b\le s\le t\le d.
\]
If $\sigma(z)\ge\sigma(w)$, this implies
\begin{align*}
 m|\sigma(z)-\sigma(w)|
 &\le U(z,\sigma(z))-U(z,\sigma(w))\\
 &=U(w,\sigma(w))-U(z,\sigma(w))
 \le C|z-w|^\alpha.
\end{align*}
Interchanging $z,w$ gives the same bound in the other case. Hence
\begin{equation}\label{eq:holder-radius}
 |\sigma(z)-\sigma(w)|\le(C/m)|z-w|^\alpha,
 \qquad 0<\alpha\le1.
\end{equation}

Finally, suppose $U\in C^{k,\alpha}_{\mathrm{loc}}$ with $k\ge1$.
The implicit function theorem first gives $\sigma\in C^k$. On a smaller
convex coordinate neighborhood, $x\mapsto(x,\sigma(x))$ is
Lipschitz and $U_t(x,\sigma(x))$ is bounded below by a positive
constant. Differentiating the level equation gives
\[
 \sigma_{x_j}=-\frac{U_{x_j}(x,\sigma(x))}{U_t(x,\sigma(x))},
 \qquad
 U_t(x,\sigma(x))D^\beta\sigma=-P_\beta,
 \quad |\beta|=k,
\]
where $P_\beta$ is a finite sum of products of derivatives of $U$
of order at most $k$, evaluated on the graph, and derivatives of
$\sigma$ of order at most $k-1$. For $k=1$, the first formula
has $\alpha$-H\"older numerator and denominator. For $k\ge2$,
\[
 D^j\sigma\in C^{0,1}_{\mathrm{loc}}\ (j<k),\qquad
 D^jU\in C^{0,1}_{\mathrm{loc}}\ (j<k),\qquad
 D^kU\in C^{0,\alpha}_{\mathrm{loc}}.
\]
Composition with the Lipschitz graph and multiplication preserve
these local H\"older estimates. Division by $U_t$ therefore gives
$D^\beta\sigma\in C^{0,\alpha}_{\mathrm{loc}}$. Formula
\eqref{eq:radius-metric} now proves the
assertion in every class.
\end{proof}

\begin{samepage}
\begin{theorem}[Regularity criterion]\label{cor:criterion}\label{thm:regularity-criterion}
Let $A\to Z$ be a holomorphic line bundle over a compact complex
manifold. For each class $\mathcal R$ in
\eqref{eq:regularity-classes}, the following are equivalent.
\begin{enumerate}[label=\textup{(\roman*)},leftmargin=2.5em]
\item $\Tot(A^{-1})$ admits a psh exhaustion of class $\mathcal R$.
\item $A$ admits a semipositive Hermitian metric of class $\mathcal R$.
\end{enumerate}
\end{theorem}
\begin{proof}
Proposition~\ref{prop:radius} gives \textup{(i)}$\Rightarrow$\textup{(ii)}.
Lemma~\ref{lem:norm} gives \textup{(ii)}$\Rightarrow$\textup{(i)}.
\end{proof}
\end{samepage}

For $C^2$ data, the relation between the two Levi forms is explicit.
Differentiating $U(z,\sigma(z))=c$ twice and using
$V(z,\zeta)=U(z,\operatorname{Re}\zeta)$ gives
\begin{equation}\label{eq:radius-Levi}
 \Levi_{-\sigma}(\xi)
 =\frac{\Levi_V(\xi,2\partial\sigma(\xi))}{U_t}\ge0.
\end{equation}
The identity is evaluated along the graph $t=\sigma(z)$.
For the dual norm, direct differentiation and completion of the
square give
\[
 \Levi_{|\lambda|^2e^\varphi}(\xi,\eta)
 =e^\varphi\left(
    |\eta+\lambda\partial\varphi(\xi)|^2
    +|\lambda|^2\Levi_\varphi(\xi)\right)\ge0.
\]

\begin{proof}[Proof of Theorem~\ref{thm:main}]\label{sec:conclusion}
Take the ruled surface $S$, the section $Y$, and the line bundle $L$
constructed in Section~\ref{sec:metrics}, for example with $F=K_B$.
The total space $M=\Tot(L^{-1})$ has the compactification
\begin{equation}\label{eq:compactification}
 M=\PP(\OO_S\oplus L^{-1})
 \setminus\PP(L^{-1}),\qquad
 v\longmapsto[(1,v)].
\end{equation}
The inverse sends $[(a,b)]$, with $a\ne0$, to $b/a$.
Since $S$ is projective and $L$ is algebraic, the projective bundle
in \eqref{eq:compactification} is projective. Therefore
$\dim_\C M=3$, $M$ is smooth and quasi-projective, and the
restriction of a projective K\"ahler form makes it K\"ahler.

Let $h$ be the locally Lipschitz semipositive metric of
Proposition~\ref{prop:continuous-metric}. Lemma~\ref{lem:norm}
gives the locally Lipschitz psh exhaustion
\begin{equation}\label{eq:explicit-exhaustion}
 \rho(v)=\|v\|_{h^{-1}}^2.
\end{equation}
If $M$ admitted a $C^1$ psh exhaustion,
Theorem~\ref{cor:criterion} would give a $C^1$ semipositive
Hermitian metric on $L$, contrary to Proposition~\ref{prop:no-metric}.
\end{proof}

\section{A non-K\"ahler modification over one point}\label{sec:nonkahler}

We first show that a modification with a compact center preserves
the existence of a $C^r$ psh exhaustion. We then place Hironaka's
opposite-order blowup construction inside the exceptional plane over
a point to prove Theorem~\ref{thm:nonkahler-threefold}.

\begin{proposition}\label{prop:compact-modification}
Let $f:X'\to X$ be a proper surjective holomorphic map of complex
manifolds, and suppose that
$f:X'\setminus f^{-1}(K)\to X\setminus K$ is a biholomorphism for
some compact $K\subset X$. For $r\in\mathbb Z_{\ge0}\cup\{\infty\}$,
$X'$ admits a $C^r$ psh
exhaustion if and only if $X$ does.
\end{proposition}
\begin{proof}
If $u$ is a $C^r$ psh exhaustion of $X$, then $u\circ f$ is
$C^r$ and psh, and
\[
 \{u\circ f\le a\}=f^{-1}(\{u\le a\})
\]
is compact by properness.

Conversely, let $F$ be a $C^r$ psh exhaustion of $X'$.
Choose a relatively compact open neighborhood $U$ of $K$ and take
$c>\sup_{f^{-1}(\overline U)}F$. Let $\chi$ be a smooth convex
nondecreasing function, zero on $(-\infty,c]$ and tending to
$+\infty$ as $t\to+\infty$.
Let $j:X\setminus K\to X'\setminus f^{-1}(K)$ be the
holomorphic inverse of $f$. Then
\[
 G=\chi\circ F\circ j\quad\text{on }X\setminus K,
 \qquad G=0\quad\text{on }K
\]
vanishes throughout $U$, because $F<c$ on $f^{-1}(U)$.
It is therefore $C^r$ across $K$. Convex nondecreasing composition
gives plurisubharmonicity on $X\setminus K$, and $G$ is constant on $U$.
The identity $G\circ f=\chi\circ F$ gives
\[
 \{G\le a\}=f(\{\chi\circ F\le a\}).
\]
The right-hand side is compact, since $\chi\circ F$ is an exhaustion.
\end{proof}

\begin{proposition}\label{prop:point-modification}
For every point $x$ of a connected smooth complex threefold $X$,
there is a connected smooth non-K\"ahler threefold $N$ and a proper
modification $\mu:N\to X$ that is biholomorphic over
$X\setminus\{x\}$.
\end{proposition}
\begin{proof}
Following the construction in \cite[\S6]{Barlet2015}, first blow up $x$,
\[
 b:W=\Bl_xX\to X,\qquad E=b^{-1}(x)\simeq\PP^2.
\]
In homogeneous coordinates $[u:v:w]$ on $E$, the line and conic
\[
 A=\{w=0\},\qquad B=\{uv=w^2\}
\]
meet transversely in $E$ at $p=[1:0:0]$ and $q=[0:1:0]$.
Set $W_p=W\setminus\{q\}$ and $W_q=W\setminus\{p\}$.
On $W_p$ and $W_q$ respectively, form
\[
 \begin{split}
 N_p&=\Bl_{\widetilde{B\cap W_p}}
             \bigl(\Bl_{A\cap W_p}W_p\bigr),\\
 N_q&=\Bl_{\widetilde{A\cap W_q}}
             \bigl(\Bl_{B\cap W_q}W_q\bigr).
 \end{split}
\]
Here a tilde denotes strict transform. Over
$W_p\cap W_q=W\setminus\{p,q\}$ the two centers are disjoint.
The two blowup orders have a canonical isomorphism there, agreeing
with the identity off $A\cup B$. Glue by this isomorphism to
obtain $N$, with maps
\[
 \nu:N\longrightarrow W,\qquad \mu=b\circ\nu:N\longrightarrow X.
\]
The centers are smooth, so $N$ is smooth. Points with distinct images
in $W$ can be separated using neighborhoods in $W$, while points over
the same image lie together in $N_p$ or $N_q$. Thus $N$ is Hausdorff.
Properness is local on the base, so $\nu$ and $\mu$ are proper.
The two connected blowup charts overlap, making $N$ connected.
Every center lies in $E=b^{-1}(x)$, so $\mu$ is biholomorphic over
$X\setminus\{x\}$.

The obstruction to a K\"ahler form is a homology relation between
two exceptional curves, as in \cite[Tag~08J1]{StacksHironaka}.
Let $D_A,D_B$ be the exceptional surfaces dominating $A,B$, with
the induced maps
\[
 f_A:D_A\longrightarrow A,\qquad f_B:D_B\longrightarrow B.
\]
Locally at a crossing they are ruled surfaces or blowups of ruled
surfaces at a point. They are therefore smooth, and properness of
$\nu$ makes them compact.
Let $\alpha,\beta\in H_2(N,\R)$ be their general fiber classes.
Since $A$ and $B$ are rational, their point divisors are linearly
equivalent. Pulling back shows that all fibers of $f_A$ represent
$\alpha$, and all fibers of $f_B$ represent $\beta$.

Near $p$, choose coordinates $(z_1,z_2,z_3)$ with
$A=\{z_2=z_3=0\}$ and $B=\{z_1=z_3=0\}$.
The first blowup has the chart
\[
 (z_1,z_2,t)\longmapsto(z_1,z_2,z_3=z_2t).
\]
Its exceptional divisor is $z_2=0$, and the strict transform of
$B$ is $z_1=t=0$. A chart of the second blowup is
\[
 (r,z_2,t)\longmapsto(z_1=rt,z_2,t).
\]
On the strict transform $D_A=\{z_2=0\}$, the map $f_A$ is
$(r,t)\mapsto z_1=rt$. Its fiber over $p$ is therefore
\[
 \operatorname{div}(rt)=\{r=0\}+\{t=0\}=R_p+T_p.
\]
The closures of these two components are respectively the strict
transform of the first exceptional fiber and the exceptional curve
of its surface blowup. Both are smooth rational curves, and the
factorization $rt$ gives multiplicity one. The second exceptional
surface $D_B$ has fiber $T_p$ over $p$. Interchanging $A$ and $B$
in the same calculation at $q$ gives
\[
 \alpha=[R_p]+[T_p],\quad\beta=[T_p],\qquad
 \beta=[R_q]+[T_q],\quad\alpha=[T_q].
\]
Subtracting each pair gives
$[R_p]=\alpha-\beta$ and $[R_q]=\beta-\alpha$, hence
\begin{equation}\label{eq:hironaka-zero-cycle}
 [R_p]+[R_q]=0\quad\text{in }H_2(N,\R).
\end{equation}
Integration of a K\"ahler form over this homologically trivial
effective cycle would give
\[
 0=\int_{R_p}\omega+\int_{R_q}\omega>0,
\]
a contradiction.
\end{proof}

\begin{proof}[Proof of Theorem~\ref{thm:nonkahler-threefold}]
Apply Proposition~\ref{prop:point-modification} to $M$ at $x$, and let
$\rho$ be the $C^{0,1}_{\mathrm{loc}}$ psh exhaustion from
Theorem~\ref{thm:main}. Since $\mu$ is smooth and proper,
$\rho\circ\mu$ is a $C^{0,1}_{\mathrm{loc}}$ psh exhaustion.
A $C^{1,0}$ psh exhaustion on
$N$ would descend after truncation by
Proposition~\ref{prop:compact-modification}, contradicting
Theorem~\ref{thm:main}.
\end{proof}

\section{Surface regularization and boundary criteria}\label{sec:surfaces}
Throughout this section, $X=S\setminus\supp D$, where $S$ is a smooth
compact connected complex surface and $D\ne0$ is effective. We identify
reduced divisors with their supports. A \emph{divergent psh collar}
along a boundary component $E$ is a psh function on $U\setminus E$,
for a neighborhood $U$ of $E$, tending uniformly to $+\infty$ as one
approaches $E$. We prove regularization by constructing smooth collars
and extending them to an exhaustion of $X$.

\subsection{Boundary tools}\label{sec:surface-pair-reduction}

The intersection matrix of a boundary component separates three cases.
Negative definiteness obstructs an exhaustion, a positive eigenvalue
gives a strictly psh collar, and a zero largest eigenvalue gives a
divisor of square zero.

\begin{lemma}\label{lem:surface-boundary-tools}
\begin{enumerate}[label=\textup{(\roman*)},leftmargin=2.5em]
\item\label{item:negative-boundary}
If a connected component of $\supp D$ has negative definite
intersection matrix, $X$ has no continuous psh exhaustion.
\item\label{item:collar-gluing}
Smooth divergent psh collars at all connected boundary components
give a smooth psh exhaustion of $X$.
\item\label{item:matrix-alternatives}
For a connected reduced divisor $E=\bigcup E_i$ with simple normal
crossings (SNC), let $\lambda$
be the largest eigenvalue of $I=(E_i\cdot E_j)$. If $\lambda>0$,
$E$ has a smooth strictly psh divergent collar. If $\lambda=0$,
$I\le0$, $\dim\ker I=1$, and
\begin{equation}\label{eq:surface-null-boundary}
 A=\sum_i a_iE_i,\qquad a_i\in\mathbb Z_{>0},\qquad
 A\cdot E_i=A^2=0.
\end{equation}
When $S$ is projective, $A$ is nef.
\end{enumerate}
\end{lemma}
\begin{proof}
For \textup{(i)}, let $E$ be the boundary component. Grauert's theorem
\cite[\S4, subsection~8(e)]{Grauert1962} contracts $E$ to
a point $p$ of a normal surface $Z$. Write the contraction as
$q:W\to Z$, with $W\setminus E\to Z\setminus\{p\}$ a biholomorphism.
An exhaustion $\rho$ on $X$ would therefore give
$u=\rho\circ(q|_{W\setminus E})^{-1}$, psh and divergent at $p$.
Choose a finite local projection
\[
 \pi:V\longrightarrow\Delta^2,\qquad \pi^{-1}(0)=\{p\},
 \qquad v(z)=\max_{y\in\pi^{-1}(z)}u(y).
\]
The function $v$ is continuous on $\Delta^2\setminus\{0\}$.
Off the branch locus it is a finite maximum of psh functions.
Properness of $\pi$ gives local boundedness at each remaining point
away from $0$, so removability makes $v$ psh there as well.
Fix $0<r<1$ and set
\[
 M_r=\max_{\{|z_1|\le r,\ |z_2|=r\}}v<\infty.
\]
For $0<|z_1|<r$, the subharmonic maximum principle on the vertical
disc gives $v(z_1,z_2)\le M_r$ for $|z_2|<r$. Continuity gives the
same bound when $z_1=0$ and $z_2\ne0$. Hence $u\le M_r$ near $p$,
contradicting its divergence.

For \textup{(ii)}, choose disjoint neighborhoods $W_E$ of the
boundary components so that each collar $\Phi_E$ is defined on a
neighborhood of $\overline W_E\setminus E$.
Choose $T_E>\sup_{\partial W_E}\Phi_E$ and a smooth convex
nondecreasing $\chi_E$, zero on $(-\infty,T_E]$ and tending to
$+\infty$. Define
\[
 \rho_E(x)=
 \begin{cases}\chi_E(\Phi_E(x)),&x\in W_E\setminus E,\\
 0,&x\notin W_E.
 \end{cases}
\]
The two formulas agree on a neighborhood of $\partial W_E$, and
\[
 \Levi_{\chi_E\circ\Phi_E}(\xi)
 =\chi_E'(\Phi_E)\Levi_{\Phi_E}(\xi)
  +\chi_E''(\Phi_E)|\partial\Phi_E(\xi)|^2\ge0.
\]
Thus $\rho=\sum_E\rho_E$ is smooth and psh. Every sequence leaving
compact subsets of $X$ approaches $D$ along a subsequence, where the
corresponding summand tends to infinity. This proves properness.

For \textup{(iii)}, Perron--Frobenius applied to $I+c\mathrm{Id}$
shows that $\lambda$ is simple and has a strictly positive
eigenvector $v$. When $\lambda>0$, both $v>0$ and $Iv>0$ persist
under small perturbations. Approximate $v$ by a rational vector and
clear denominators to obtain
$A=\sum a_iE_i$ with $a_i\in\mathbb Z_{>0}$ and
$A\cdot E_i>0$ for all $i$.
Thus $\OO_S(A)|_E$ is ample. A projective embedding by a power
gives a Fubini--Study metric on this restriction. Extend its weights
locally to $S$ and patch them relative to a reference metric.
The extensions have positive tangential curvature along $E$.
Near a node, lift the embedding sections from $E$ to a coordinate
neighborhood in $S$. Their ratios have differentials spanning the
two-dimensional Zariski cotangent space of the node, and hence the
ambient cotangent space. Their Fubini--Study weights have positive
curvature in both directions. Keep these weights on smaller node
neighborhoods when patching. The other extensions agree along $E$,
so the patched metric retains positive tangential curvature at its
smooth points.

Let $q=\|s_E\|^2$ for a smooth metric on $\OO_S(E)$.
At a smooth point of $E=\{w=0\}$, write $q=|w|^2e^{-\gamma}$.
Then
\[
 \partial\barp q|_E
   =e^{-\gamma}|_E\,dw\wedge d\bar w.
\]
Away from the node neighborhoods the curvature matrix of the
patched metric has the form
$\left(\begin{smallmatrix}a&b\\\bar b&d\end{smallmatrix}\right)$
in tangential and normal directions, with $a\ge\varepsilon>0$.
Adding $Cq$ to its weights replaces $d$ by $d+Ce^{-\gamma}$ along $E$.
By compactness, one constant $C$ makes
$a(d+Ce^{-\gamma})>|b|^2$. At the nodes $\ddc q=0$, and at other
points of $E$ it is semipositive, so positivity is preserved there.
Continuity now gives a metric $h$ of positive curvature near all
of $E$. On the complement,
\[
 \ddc\bigl(-\log\|s_A\|_h^2\bigr)=\Theta_h(\OO_S(A))>0,
\]
and $a_i>0$ gives uniform divergence. For $\lambda=0$, every other
eigenvalue is negative. Since $I$ has integer entries and rank one
less than its size, its positive kernel has a rational generator, and
clearing denominators proves \eqref{eq:surface-null-boundary}.
Effectiveness gives $A\cdot C\ge0$ for every curve $C\not\subset E$,
which proves nefness in the projective case.
\end{proof}

For a boundary resolution $p:\widetilde S\to S$, the restriction
$\widetilde S\setminus p^{-1}(\supp D)\to X$ is a biholomorphism.
Interior blowups preserve both exhaustion properties by
Proposition~\ref{prop:compact-modification}.

\begin{proposition}\label{prop:surface-nonconstant-function}
A connected noncompact smooth complex surface $X$ with a continuous psh
exhaustion and $\OO(X)\ne\C$ has a smooth psh exhaustion.
\end{proposition}
\begin{proof}
V\^{\i}j\^{\i}itu's theorem \cite[Theorem~1.1]{Vajaitu2022} gives
holomorphic convexity. Let $r:X\to Y$ be the proper Remmert
reduction and let $\tau$ be a smooth strictly psh exhaustion of the
Stein space $Y$. Then
\[
 \rho=\tau\circ r,\qquad
 \{\rho\le c\}=r^{-1}(\{\tau\le c\})\Subset X
\]
is a smooth psh exhaustion.
\end{proof}

\begin{proposition}\label{cor:curve-total-space}
Let $A$ be a line bundle on a compact Riemann surface $B$.
The following are equivalent.
\begin{enumerate}[label=\textup{(\roman*)},leftmargin=2.5em]
\item $\Tot(A^{-1})$ admits a continuous psh exhaustion.
\item $\deg A\ge0$.
\item $\Tot(A^{-1})$ admits a smooth psh exhaustion.
\end{enumerate}
\end{proposition}
\begin{proof}
For \textup{(i)}$\Rightarrow$\textup{(ii)},
Theorem~\ref{cor:criterion} gives a continuous semipositive metric on $A$.
Its positive curvature current has total mass $\deg A$, so $\deg A\ge0$.

For \textup{(ii)}$\Rightarrow$\textup{(iii)}, put $d=\deg A\ge0$.
Choose an area form $\omega$ of integral one and a smooth reference
metric $h_0$ of curvature $\theta_0$. Since $d\omega-\theta_0$ has
integral zero, the equation $\ddc f=d\omega-\theta_0$ has a smooth
solution on $B$. The metric $h=e^{-f}h_0$ then satisfies
\[
 \Theta_h(A)=\theta_0+\ddc f=d\omega\ge0.
\]
Lemma~\ref{lem:norm} gives the smooth psh exhaustion.
The implication \textup{(iii)}$\Rightarrow$\textup{(i)} is immediate.
\end{proof}

\subsection{Current averaging and regularization}
\label{subsec:divisor-regularization}

The main step is to average currents supported on exhaustion levels.
The resulting class is orthogonal to the boundary, and the Hodge index
theorem relates it to a divisor of square zero.
Theorem~\ref{thm:surface-introduction} follows from the next theorem
and its quasi-projective corollary.

\begin{theorem}\label{thm:divisor-regularization}
Let $S$ be a compact K\"ahler surface and $D=\sum_iD_i\ne0$ a reduced
divisor. Suppose $X=S\setminus D$ admits a continuous psh exhaustion.
Then it admits a smooth psh exhaustion. If $X$ is not a modification
of a Stein space, there is a divisor
\[
 A=\sum_i a_iD_i,\qquad a_i\in\mathbb Z_{>0},\qquad
 A\cdot D_i=A^2=0,
\]
such that $\OO_S(A)$ has a smooth Hermitian metric with zero
curvature near $D$. Its canonical section gives a pluriharmonic
divergent collar along $D$.
\end{theorem}

\begin{lemma}\label{lem:level-averaging}
Let $(S,\omega)$ be a compact K\"ahler surface and $X\subset S$
a connected noncompact open set with a continuous psh exhaustion
$\rho$. If $X$ is not a modification of a Stein space, there is
a positive $(1,1)$-current $T$ on $S$ such that
\begin{equation}\label{eq:averaged-current-properties}
 \ddc T=0,\qquad \int_S T\wedge\omega=1,\qquad
 \supp T\Subset X,\qquad \nu(T,x)=0\quad(x\in S).
\end{equation}
\end{lemma}
\begin{proof}
Let $\Sigma_X^0$ be the set of points at which every continuous
psh exhaustion fails to be strongly psh. By
\cite[Lemma~2.1]{MongodiSlodkowski2020},
\begin{equation}\label{eq:kernel-halfline}
 (t_0,\infty)\subset\rho(\Sigma_X^0)
\end{equation}
for some $t_0$. For $t>t_0$, no smooth strictly psh function exists
near $K_t=\{\rho=t\}$. Indeed, if $v$ were such a function,
properness permits $\delta>0$ with the compact shell
$\{t-\delta\le\rho\le t+\delta\}$ contained in its domain.
After scaling and adding a constant, arrange $0<v<1$ there. For
$0<\eta<1$, $0<\varepsilon<\eta\delta/2$, set
$g(s)=s+\eta|s-t|$. The function
\[
 \psi=
 \begin{cases}
 \max\{g(\rho),\rho+\varepsilon v\},& |\rho-t|<\delta,\\
 g(\rho),& |\rho-t|\ge\delta
 \end{cases}
\]
is well defined because the formulas agree near $|\rho-t|=\delta$,
where $g(\rho)-\rho=\eta|\rho-t|>\varepsilon v$. Both entries of the
maximum are psh, since $g$ is convex and increasing.
On $K_t$ the second entry strictly dominates, so
$\psi=\rho+\varepsilon v$ is strongly psh near $K_t$.
Finally $\psi\ge g(\rho)$ makes $\psi$ an exhaustion. This
contradicts $K_t\cap\Sigma_X^0\ne\varnothing$.

Let $\mathcal B_t$ be the Banach space of continuous real
$(1,1)$-forms on $K_t$, with the uniform norm induced by $\omega$.
Consider the subspace and open cone
\[
 V_t=\{\ddc f|_{K_t}:f\in C^\infty(S,\R)\},\qquad
 P_t=\{\eta\in\mathcal B_t:\eta>0\text{ on }K_t\}.
\]
The absence of a smooth strictly psh function near $K_t$ gives
$V_t\cap P_t=\varnothing$. Hahn--Banach separation therefore
gives a nonzero continuous functional $\ell_t$ with
\[
 \ell_t|_{V_t}=0,\qquad \ell_t(\eta)\ge0\quad(\eta\ge0).
\]
The second assertion follows by applying the separating inequality
to $\eta+\varepsilon\omega$ and letting $\varepsilon\downarrow0$.
Moreover $\ell_t(\omega)>0$: otherwise
$-C\omega\le\eta\le C\omega$ would force $\ell_t(\eta)=0$ for
every $\eta$. Define the current by
\[
 \langle T_t,\eta\rangle
   =\frac{\ell_t(\eta|_{K_t})}{\ell_t(\omega|_{K_t})}.
\]
This defines a positive current supported on $K_t$, and
$\ell_t|_{V_t}=0$ gives $\ddc T_t=0$. Hence
\begin{equation}\label{eq:level-current}
 T_t\ge0,\qquad \ddc T_t=0,\qquad
 \supp T_t\subset K_t,\qquad \int_S T_t\wedge\omega=1.
\end{equation}

To average these currents, fix $t_0<a<b$ and set
\[
 t_{n,j}=a+\frac{j-1/2}{n}(b-a),\qquad
 T_n=\frac1n\sum_{j=1}^nT_{t_{n,j}}.
\]
Each $T_n$ has mass one and support in the compact shell
$K=\rho^{-1}([a,b])\Subset X$. A weakly convergent subsequence has
a positive limit $T$ satisfying
\[
 \int_S T\wedge\omega=1,\qquad \ddc T=0,
 \qquad\supp T\subset K.
\]

Fix a coordinate ball $B(x,2R)\Subset X$, and put
$\beta=\ddc|z|^2$. Lelong--Jensen monotonicity for positive
$\ddc$-closed currents \cite[Appendix~B, (B.1)--(B.2)]{DNS22} gives,
uniformly in $t$, for $0<r<R$,
\begin{equation}\label{eq:uniform-lelong-bound}
 r^{-2}\int_{B(x,r)}T_t\wedge\beta
 \le R^{-2}\int_{B(x,R)}T_t\wedge\beta\le C_R.
\end{equation}
Indeed, $\beta\le c_R\omega$ on $B(x,R)$, so one may take
$C_R=c_R/R^2$, independently of $t$. Put
$m_r=\inf_{B(x,r)}\rho$ and $M_r=\sup_{B(x,r)}\rho$.
Since $\supp T_t\subset\{\rho=t\}$, only indices with
$t_{n,j}\in[m_r,M_r]$ contribute. Their number $N_n(r)$ satisfies
\[
 N_n(r)\le\frac{n(M_r-m_r)}{b-a}+2.
\]
Summing \eqref{eq:uniform-lelong-bound} over these indices gives
\begin{equation}\label{eq:grid-oscillation-bound}
 r^{-2}\int_{B(x,r)}T_n\wedge\beta
 \le C_R\left(\frac{\operatorname{osc}_{B(x,r)}\rho}{b-a}
                  +\frac2n\right).
\end{equation}
Multiply $\beta$ by a cutoff equal to one on $B(x,R)$. The
corresponding positive measures $\mu_n$ converge weakly to $\mu$.
For the open ball, the Portmanteau inequality and
\eqref{eq:grid-oscillation-bound} give
\[
 \mu(B(x,r))\le\liminf_n\mu_n(B(x,r))
 \le\frac{C_Rr^2}{b-a}(M_r-m_r).
\]
Consequently,
\begin{equation}\label{eq:limit-oscillation-bound}
 r^{-2}\int_{B(x,r)}T\wedge\beta
 \le\frac{C_R}{b-a}\operatorname{osc}_{B(x,r)}\rho
 \longrightarrow0\qquad(r\downarrow0).
\end{equation}
Thus $\nu(T,x)=0$ on $X$. Since $\supp T\Subset X$, the Lelong
numbers vanish on all of $S$.
\end{proof}

\begin{lemma}\label{lem:surface-current-inequality}
For a positive $\ddc$-closed $(1,1)$-current $T$ on a compact
K\"ahler surface with all Lelong numbers zero, there is a class
$\alpha_T\in H^{1,1}(S,\R)$ satisfying
\begin{equation}\label{eq:current-class-pairing}
 \alpha_T\cdot[\eta]=\int_S T\wedge\eta
 \qquad(d\eta=0),
\end{equation}
for every smooth real $(1,1)$-form $\eta$. Moreover,
\begin{equation}\label{eq:DNS-inequality}
 \alpha_T^2\ge0,\qquad dT\ne0\ \Longrightarrow\ \alpha_T^2>0.
\end{equation}
If $dT=0$, then $\alpha_T=[T]$.
\end{lemma}
\begin{proof}
The decomposition \cite[\S2.1 and Lemma~B.2]{DNS22}
\[
 T=\Omega+\partial R+\barp\overline R
                  +i\partial\barp u
\]
has $\Omega$ smooth, real, and closed, $R$ of bidegree $(0,1)$,
and $u$ a real distribution. For $d\eta=0$, its pure bidegree
implies $\partial\eta=\barp\eta=0$. Thus
\[
 \langle\partial R,\eta\rangle
 =\langle\barp\overline R,\eta\rangle
 =\langle i\partial\barp u,\eta\rangle=0,
\]
which proves \eqref{eq:current-class-pairing} with
$\alpha_T=[\Omega]$. The nondegenerate intersection pairing on
$H^{1,1}(S,\R)$ makes the class independent of the decomposition.
The density formula \cite[Theorem~2.2(2), Corollary~2.4]{DNS22} gives
\[
 \alpha_T^2=\|\vartheta\|+2E(T),\qquad
 \|\vartheta\|\ge0,\quad E(T)\ge0,
 \quad E(T)=0\Longleftrightarrow dT=0.
\]
Here $\vartheta$ is the positive measure on the diagonal determined
by a tangent current of $T\otimes T$, and
$E(T)=\int_S\barp R\wedge\partial\overline R$ is the
current energy, independent of the chosen decomposition.
The hypotheses of the density formula hold because all Lelong numbers vanish. This
proves \eqref{eq:DNS-inequality}.
\end{proof}

\begin{proof}[Proof of Theorem~\ref{thm:divisor-regularization}]
If $r:X\to Y$ is a proper modification of a Stein space, the
pullback of a smooth strictly psh exhaustion of $Y$ is the required
exhaustion. Otherwise the preceding lemmas give a current $T$ with
support in $K\Subset X$ and a class $\alpha=\alpha_T$ satisfying
\begin{equation}\label{eq:alpha-basic}
 \alpha\cdot[\omega]=1,\qquad \alpha^2\ge0,
 \qquad dT\ne0\Longrightarrow\alpha^2>0.
\end{equation}
For each component $D_i$, choose a metric of curvature $\theta_i$
and canonical section $s_i$. Choose $\chi\in C_c^\infty(X)$ equal
to one near $K$, and extend
$v_i=\chi\log\|s_i\|^2$ by zero to $S$.
Since $s_i$ is nonvanishing near $K$, Poincar\'e--Lelong gives
$\ddc v_i=-\theta_i$ there. Therefore
\begin{equation}\label{eq:alpha-boundary-orthogonality}
 \alpha\cdot[D_i]=\int_S T\wedge\theta_i
          =-\langle T,\ddc v_i\rangle=0.
\end{equation}

For a connected boundary component $E$, let $\lambda_E$ be the
largest eigenvalue of its intersection matrix $I_E$.
Lemma~\ref{lem:surface-boundary-tools}\textup{(i)} excludes
$\lambda_E<0$. If $\lambda_E>0$, rational approximation gives an
integral divisor $B=\sum b_iD_i$, with all $b_i>0$ and
$B\cdot D_i>0$ on $E$. In particular,
$B^2=\sum b_i(B\cdot D_i)>0$.
Then $\alpha\cdot[B]=0$, and Hodge index forces $\alpha^2<0$,
contrary to \eqref{eq:alpha-basic}. Therefore $\lambda_E=0$.
Distinct irreducible curves have nonnegative intersection, and connectedness
makes $I_E$ irreducible. Perron--Frobenius gives a one-dimensional
kernel with a positive rational generator. Choosing a positive integral
generator for each $E$ and summing gives
\begin{equation}\label{eq:full-support-null-divisor}
 A=\sum_i a_iD_i,\qquad a_i>0,\qquad
 \supp A=D,\qquad A\cdot D_i=A^2=0.
\end{equation}

Set $c=(\int_A\omega)^{-1}$ and $v=\alpha-c[A]$.
Effectiveness gives $\int_A\omega>0$, while
\eqref{eq:alpha-basic}--\eqref{eq:alpha-boundary-orthogonality} give
\[
 v\cdot[\omega]=0,\qquad v\cdot[A]=0,\qquad v^2=\alpha^2\ge0.
\]
Hodge index is negative definite on $[\omega]^\perp$, so $v=0$
and $\alpha^2=0$. The strict inequality in
\eqref{eq:alpha-basic} then forces $dT=0$. We have proved
\begin{equation}\label{eq:closed-current-null-class}
 dT=0,\qquad [T]=\alpha=c[A],\qquad
 c=\left(\int_A\omega\right)^{-1}>0.
\end{equation}
The $\ddc$-lemma for currents now supplies a real distribution $u$ with
\begin{equation}\label{eq:global-log-potential}
 \ddc u=T-c[A].
\end{equation}
Choose a neighborhood $U$ of $D$ disjoint from $K$.
On a trivializing chart $U_i\subset U$, write
$s_A=f_i e_i$. Poincar\'e--Lelong and
\eqref{eq:global-log-potential} give the distributional identity
\[
 \ddc\bigl(u+c\log|f_i|^2\bigr)
 =-c[A]+c[A]=0.
\]
Elliptic regularity for the pluriharmonic distribution
$u+c\log|f_i|^2$ gives
\begin{equation}\label{eq:local-log-collar}
 u=-c\log|f_i|^2+h_i,
 \qquad h_i\text{ smooth and pluriharmonic on }U_i.
\end{equation}
The functions $-c\log|f_i|^2+h_i$
agree on overlaps off $D$, since they represent the same distribution $u$.
They therefore define a smooth pluriharmonic function on
$U\setminus D$. In a finite relatively compact chart cover of $D$,
each $h_i$ is bounded and $f_i$ tends to zero along $D$. Since
every $a_i>0$, the function tends uniformly to $+\infty$ along $D$.
Lemma~\ref{lem:surface-boundary-tools}\textup{(ii)} gives the exhaustion.

For the metric statement, set
$b_i=u/c+\log|f_i|^2=h_i/c$ on $U_i$.
If $e_j=g_{ij}e_i$, then $f_j=g_{ij}^{-1}f_i$ and
\[
 b_j=b_i-\log|g_{ij}|^2,\qquad \ddc b_i=0.
\]
Thus $\|e_i\|_{h_{\mathrm{flat}}}^2=e^{-b_i}$ defines a smooth flat
Hermitian metric on $\OO_S(A)|_U$, with
$-\log\|s_A\|_{h_{\mathrm{flat}}}^2=u/c$.
For a global smooth metric $h_0$ with weights $b_i^0$, the
differences $b_i-b_i^0$ define a single smooth function $b$ on $U$.
Choose a cutoff $\chi$ supported in $U$ and equal to one on a
smaller neighborhood of $D$. The global metric
$h=e^{-\chi b}h_0$ agrees with $h_{\mathrm{flat}}$ there, proving the
claimed curvature statement.
\end{proof}

\begin{corollary}\label{cor:quasiprojective-regularization}
Every smooth complex quasi-projective surface with a continuous
psh exhaustion has a smooth psh exhaustion.
\end{corollary}
\begin{proof}
Resolve a projective closure and principalize the complement ideal,
with centers outside $X$ \cite[Theorems~1.0.1--1.0.2]{Wlodarczyk05}.
The resulting morphism $p:S\to\overline X$ is an isomorphism over
$X$, and $D=(p^{-1}(\overline X\setminus X))_{\mathrm{red}}$ is SNC.
Thus $p$ identifies $S\setminus D$ with $X$. Apply
Theorem~\ref{thm:divisor-regularization} to each noncompact component.
Compact components need only constant functions, and there are
finitely many components.
\end{proof}

For an example with $\OO(X)=\C$, take a compact Riemann surface
$B$ of positive genus and a non-torsion $A\in\operatorname{Pic}^0(B)$.
The squared unitary-flat norm is a smooth psh exhaustion of
$X=\Tot(A^{-1})$. The degree-$n$ fiberwise Taylor coefficient of a
holomorphic function on $X$ belongs to $H^0(B,A^n)=0$ for $n>0$,
and its constant coefficient is constant on $B$.

\begin{corollary}\label{cor:square-zero-continuous-collar}
Let $C\subset S$ be a smooth connected square-zero curve in a compact
K\"ahler surface. A continuous divergent psh collar at $C$ implies
a smooth psh exhaustion of $S\setminus C$. If the complement is not
a modification of a Stein space, $\OO_S(C)$ is unitary flat near $C$,
and $C$ has a pluriharmonic divergent collar.
\end{corollary}
\begin{proof}
Let $\Phi$ be the collar, choose a neighborhood $W$ of $C$ such that
$\Phi$ is defined near $\overline W\setminus C$,
and take $M>\sup_{\partial W}\Phi$. Set
\[
 \rho(x)=
 \begin{cases}\max\{\Phi(x),M\},&x\in W\setminus C,\\
 M,&x\notin W.
 \end{cases}
\]
The formulas agree near $\partial W$, and divergence at $C$ makes
$\rho$ a continuous psh exhaustion. In
Theorem~\ref{thm:divisor-regularization}, the single-component kernel
generator can be chosen to be $1$, so $A=C$.
\end{proof}

\subsection{Numerical existence criteria}\label{sec:surface-nonprojective}

We next use the classification of compact surfaces to obtain a numerical
criterion for nonprojective surfaces, followed by the corresponding
criterion for smooth ruled surfaces.

\begin{theorem}\label{thm:nonprojective}
Let $S$ be a smooth compact connected nonprojective surface and
$D\ne0$ an effective divisor. The following conditions are equivalent.
\begin{enumerate}[label=\textup{(\roman*)},leftmargin=2.5em]
\item $S\setminus\supp D$ admits a continuous psh exhaustion.
\item $S\setminus\supp D$ admits a smooth psh exhaustion.
\item no connected boundary component has negative definite
intersection matrix.
\end{enumerate}
\end{theorem}
\begin{proof}
The implication \textup{(ii)}$\Rightarrow$\textup{(i)} is immediate, and
Lemma~\ref{lem:surface-boundary-tools}\textup{(i)} gives
\textup{(i)}$\Rightarrow$\textup{(iii)}. To prove
\textup{(iii)}$\Rightarrow$\textup{(ii)}, we distinguish the two possible
values of the algebraic dimension $a(S)$, using the classification in
\cite[\S2]{FN2005}.

If $a(S)=1$, the algebraic reduction is a
holomorphic elliptic fibration $f:S\to B$ with connected fibers,
and every curve is vertical. For a fiber $F=\sum m_iC_i$,
the identities $F\cdot C_i=0$ imply
$m_iC_i^2=-\sum_{j\ne i}m_j(C_i\cdot C_j)$. Substitution gives
\begin{equation}\label{eq:surface-fiber-intersection-identity}
 x^{\mathsf T}Ix=-\sum_{i<j}(C_i\cdot C_j)m_im_j
       \left(\frac{x_i}{m_i}-\frac{x_j}{m_j}\right)^2,
 \qquad I=(C_i\cdot C_j).
\end{equation}
The right-hand side vanishes precisely when all $x_i/m_i$ are
equal, since the fiber is connected. Extend a vector on a proper
subset of the components by zero. Equality then forces the entire
vector to vanish. Thus every proper principal submatrix is negative
definite. Condition \textup{(iii)} implies that the boundary contains
each fiber support it meets, and
$X=f^{-1}(B\setminus P)$ for a finite nonempty $P$.
For a smooth subharmonic exhaustion $\tau$ of $B\setminus P$,
$\tau\circ f$ is psh and
$\{\tau\circ f\le c\}=f^{-1}(\{\tau\le c\})$ is compact.

Suppose $a(S)=0$. By \cite[Lemma~2.1]{FN2005}, $c_1(L)^2\le0$
for every line bundle $L$. Perron--Frobenius and \textup{(iii)}
therefore give an effective $A$ of full boundary support with
$A^2=0$. Indeed, a positive largest eigenvalue would give an
effective divisor of positive square, and a negative one is excluded
by \textup{(iii)}. At a blowdown $q:T\to T'$ with exceptional curve
$E$, put $A'=q_*A$ and write $A=q^*A'+bE$. Since
$q^*A'\cdot E=0$ and $E^2=-1$,
\[
 0=A^2=(A')^2-b^2,\qquad (A')^2\le0,
\]
so $b=0$ and $A=q^*A'$. Iterating to the minimal model $p:S\to S_0$ gives
$A=p^*A_0\ne0$, $A_0^2=0$, and
$\supp A=p^{-1}(\supp A_0)$. The restriction
\[
 p:S\setminus\supp A\longrightarrow S_0\setminus\supp A_0
\]
is proper because the deleted set is the full inverse image of
$\supp A_0$. Pullback therefore preserves an exhaustion.

If $S_0$ is K\"ahler, it is a torus or K3 of algebraic dimension
zero. On a torus, translating an irreducible curve shows its square
is nonnegative, hence zero. Adjunction then gives arithmetic genus
one. A rational normalization is impossible, since any map from
$\PP^1$ lifts to $\C^2$ and is constant. The curve is therefore a
smooth elliptic curve, and a translate is an elliptic subtorus.
The quotient onto the one-dimensional torus gives a nonconstant
meromorphic function. On a K3, $h^0(\OO(-A_0))=0$ because $A_0$
is effective and nonzero, while Riemann--Roch and Serre duality give
\[
 h^0(\OO(A_0))-h^1(\OO(A_0))
 =2+\tfrac12 A_0^2=2.
\]
Two independent sections have a nonconstant meromorphic ratio.
Both possibilities contradict $a(S_0)=0$.

The remaining minimal surfaces are of class~VII. If $b_2(S_0)=n>0$, Enoki's
square-zero-divisor theorem \cite[Main Theorem, p.~454]{Enoki}
makes $S_0$ an Enoki surface and gives
\[
 A_0=mC_{\mathrm E},\qquad m\in\mathbb Z_{>0},
 \qquad \supp A_0=C_{\mathrm E},
\]
where $C_{\mathrm E}$ is its distinguished reduced cycle. By
\cite[Proposition~3.5\textup{(iii)}, p.~460]{Enoki},
$\mathcal A=S_0\setminus C_{\mathrm E}$ is an affine-line bundle
over an elliptic curve $B_{\mathrm E}$. Its linearization $Q$ has
$\deg Q=-n$, using the affine-bundle degree convention of
\cite[\S2, p.~457]{Enoki}.

Write its affine coordinate changes as
$w_i=a_{ij}w_j+b_{ij}$. Extending these maps to $\PP^1$ gives
a compact ruled surface $\overline{\mathcal A}$ with infinity
section $Y_\infty$. In the coordinates $u_i=w_i^{-1}$ near this
section,
\[
 u_i=\frac{u_j}{a_{ij}+b_{ij}u_j}
     =a_{ij}^{-1}u_j+O(u_j^2),\qquad
 N_{Y_\infty/\overline{\mathcal A}}\simeq Q^{-1}.
\]
Consequently
\[
 Y_\infty^2=\deg Q^{-1}=n>0,\qquad
 \overline{\mathcal A}\setminus Y_\infty
 \simeq S_0\setminus\supp A_0.
\]
Lemma~\ref{lem:surface-boundary-tools}\textup{(iii)} supplies a
smooth strictly psh divergent collar at $Y_\infty$, and
\textup{(ii)} gives a smooth psh exhaustion of its complement.
If $b_2(S_0)=0$, a minimal class~VII surface containing a curve is
elliptic or Hopf \cite[p.~453]{Enoki}. Since $a(S_0)=0$, it is Hopf.
Since $b_2=0$, every curve has square zero and distinct curves have
intersection zero. Adjunction gives arithmetic genus one, while
rational normalizations are impossible because maps
$\PP^1\to S_0$ lift to $\C^2\setminus\{0\}$ and are constant.
Thus the curves are disjoint smooth elliptic curves.
Finite Ueda type gives a smooth strictly psh collar
\cite[Theorem~2.5]{CLPT2018}. Infinite type gives a unitary flat
neighborhood by \cite[Theorem~1.6]{CFSTZ2026}, hence a
pluriharmonic collar. Glue the collars.
\end{proof}

\begin{theorem}\label{cor:ruled-all-degrees}
Let $f:S\to B$ be a smooth $\PP^1$-bundle over a compact Riemann
surface and let $D\ne0$ be a reduced SNC divisor on $S$.
The three conditions in Theorem~\ref{thm:nonprojective} are equivalent.
\end{theorem}
\begin{proof}
As above, it suffices to construct smooth divergent collars under
condition \textup{(iii)}. Lemma~\ref{lem:surface-boundary-tools}
handles boundary components with a positive eigenvalue. For a connected component
$E$ with zero largest eigenvalue, choose its positive kernel divisor
$A$. If $E$ is vertical, it is a whole fiber, and the punctured-base
construction applies. Otherwise, for a fiber $F$ of $f$, put
$a=A\cdot F>0$. A vertical
component would be a fiber and have intersection $a$ with $A$,
contradicting $A\cdot E_i=0$. Thus every $E_i$ is horizontal.
Choose a section $C_0$ of the projective ruling and write $C_0^2=-e$.
Write
\[
 A\equiv aC_0+bF,\qquad E_i\equiv d_iC_0+b_iF,
 \qquad d_i=E_i\cdot F>0.
\]
Using $C_0\cdot F=1$ and $F^2=0$, compute
\[
 0=A^2=-ea^2+2ab,\qquad
 0=A\cdot E_i=a b_i+d_i b-ea d_i.
\]
Since $a>0$, these identities give
\[
 b=ae/2,\qquad E_i\equiv d_i(C_0+(e/2)F),\qquad
 E_i^2=E_i\cdot E_j=0\quad(i\ne j).
\]
Distinct components are therefore disjoint, and connectedness makes
$E=C$ a single curve, smooth by the SNC hypothesis. For
$d=C\cdot F$, the canonical class formula
\[
 K_S\equiv-2C_0+(2g(B)-2-e)F
\]
and adjunction give $2g(C)-2=(K_S+C)\cdot C=d(2g(B)-2)$.
The nonconstant map $f|_C:C\to B$ is finite of degree $d$.
Riemann--Hurwitz now gives degree zero for its effective ramification
divisor, so it is \'etale.

For $d=1$, the complement $S\setminus C$ is an affine-line bundle
with degree-zero linearization. In unitary flat frames write its
transitions as $w_i=a_{ij}w_j+b_{ij}(z)$ with $|a_{ij}|=1$.
The construction of \cite[Proposition~2.1]{KoikeUeda2025} gives
pluriharmonic functions $h_i$ with
$h_i=a_{ij}h_j+b_{ij}$. On simply connected charts split
$h_i=f_i+g_i$, with $f_i$ holomorphic and $g_i$ antiholomorphic.
The coordinate changes $w_i\mapsto w_i-f_i$ replace $b_{ij}$ by
\[
 c_{ij}=b_{ij}+a_{ij}f_j-f_i=g_i-a_{ij}g_j.
\]
This expression is both holomorphic and antiholomorphic, hence
constant. In the new coordinates,
\[
 w_i=a_{ij}w_j+c_{ij},\qquad g_i=a_{ij}g_j+c_{ij},
 \qquad |a_{ij}|=1.
\]
The identity $w_i-g_i=a_{ij}(w_j-g_j)$ makes
$\rho=|w_i-g_i|^2$ global. Since $g_i$ is antiholomorphic,
\[
 \rho_{z\bar z}=|(g_i)_{\bar z}|^2,\qquad
 \rho_{z\bar w_i}=\rho_{w_i\bar z}=0,\qquad
 \rho_{w_i\bar w_i}=1.
\]
Thus $\rho$ is smooth and psh. On a finite relatively compact
trivializing cover of the compact base, the $g_i$ are bounded, so
$\rho\le c$ bounds each fiber coordinate. Its sublevels are compact.

For $d\ge2$, let $B'\to B$ be the finite \'etale Galois cover
associated to the kernel of the permutation monodromy
$\pi_1(B)\to\mathfrak S_d$ of $C\to B$. Form
\[
 S'=S\times_B B',\qquad p:S'\to S,\qquad
 C'=p^{-1}(C)=C_1\sqcup\cdots\sqcup C_d.
\]
The $C_i$ are sections with square zero because their normal bundles
are pullbacks of $N_{C/S}$ along finite covers. If $d=2$, these sections
identify the ruling with $\PP(\OO_{B'}\oplus L)$, where
$\deg L=0$, and $S'\setminus C'=L^\times$. For a unitary flat
metric, put $t(v)=\log\|v\|^2$. Then $\ddc t=0$ on $L^\times$,
and
\[
 \Levi_{t^2}(\xi)=2|\partial t(\xi)|^2\ge0,\qquad
 \{t^2\le c\}=\{e^{-\sqrt c}\le\|v\|^2\le e^{\sqrt c}\}
\]
is compact. Hence $t^2$ is a smooth psh exhaustion.

For $d\ge3$, the unique fiberwise M\"obius map sending
$C_1,C_2,C_3$ to $0,1,\infty$ identifies $S'$ with
$B'\times\PP^1$. Every remaining section is a holomorphic map
$B'\to\PP^1\setminus\{0,1,\infty\}$, which is constant by the
maximum principle. Thus $S'\setminus C'=B'\times(\PP^1\setminus Q)$
for a finite set $Q$, and an exhaustion of the punctured fiber pulls
back to one of the product.

For $d\ge2$, let $\rho'\ge0$ be the exhaustion just constructed and
$G$ the deck group of $B'\to B$. The function
\[
 \widehat\rho=\frac1{|G|}\sum_{g\in G}\rho'\circ g
\]
is smooth, psh, and $G$-invariant. Since
$\widehat\rho\ge\rho'/|G|$, its sublevels are compact.
It descends under the finite \'etale map $p$ to a smooth psh
exhaustion of $S\setminus C$. Restrict this function to a collar
at $E$ and apply Lemma~\ref{lem:surface-boundary-tools}\textup{(ii)}
to glue all boundary collars.
\end{proof}

\begin{samepage}
\subsection{Square-zero curves}\label{subsec:smooth-boundary}
For a smooth square-zero curve, Ueda theory relates the formal
neighborhood to the existence of holomorphic defining functions
$w_i=t_{ij}w_j$ with $|t_{ij}|=1$. These functions give a
pluriharmonic collar through $-\log|w_i|^2$.

\begin{proposition}\label{prop:surface-ueda-sufficient}
Let $C\subset S$ be smooth, connected, and square-zero, with $S$
compact K\"ahler. Put $N=N_{C/S}$, let
$\chi:H_1(C,\mathbb Z)\to\Uone$ be its unitary character, and set
$K=\ker[H_1(C,\mathbb Z)\to H_1(S,\mathbb Z)]$.
The complement $S\setminus C$ has a smooth psh exhaustion if any of
the following conditions holds.
\begin{enumerate}[label=\textup{(\roman*)},leftmargin=2.5em]
\item $C$ has finite Ueda type or a unitary flat neighborhood.
\item $N$ is torsion or Diophantine.
\item $\chi(K)$ is finite.
\item $q(S)\ge g(C)$.
\item $g(C)=1$.
\end{enumerate}
Here Diophantine means
$\operatorname{dist}(N^m,\OO_C)\ge c m^{-\tau}$ for some
$c,\tau>0$ and all $m\ge1$, with respect to a fixed flat metric on
$\operatorname{Pic}^0(C)$. We write
$q(S)=h^1(S,\OO_S)$ for the irregularity.
\end{proposition}
\end{samepage}
\begin{proof}
Finite type gives a smooth strictly psh collar
\cite[Theorem~2.5]{CLPT2018}. In a unitary flat neighborhood,
$w_i=t_{ij}w_j$ with $|t_{ij}|=1$, so the functions
$-\log|w_i|^2$ agree and form a pluriharmonic divergent collar.
In infinite type, Ueda's convergence theorem
\cite[Theorem~3]{Ueda1983} supplies such a neighborhood when $N$
is torsion or Diophantine.

Suppose $\chi(K)$ is finite, and let $k$ annihilate this finite
group. Write $j_*:H_1(C,\mathbb Z)\to H_1(S,\mathbb Z)$ for
the inclusion map. The formula
\[
 \chi_0(j_*\gamma)=\chi(\gamma)^k
\]
defines a character on $\operatorname{im}j_*$, because two lifts
differ by an element of $K$. Divisibility of $\Uone$ extends it
to a character $\widetilde\chi:H_1(S,\mathbb Z)\to\Uone$, with
$\widetilde\chi\circ j_*=\chi^k$. For non-torsion $N$ of infinite
type, \cite[Theorem~4.4]{CLPT2018} then gives an effective divisor
$E$ disjoint from $C$ with $[E]=m[C]$ for some $m>0$.
The line bundle $\OO_S(mC-E)$ has zero real Chern class and hence
a unitary flat metric on the compact K\"ahler surface. Its canonical
meromorphic section $\sigma=s_C^m/s_E$ satisfies
\[
 \operatorname{div}(\sigma)=mC-E,\qquad
 \ddc\log\|\sigma\|^2=m[C]-[E].
\]
On a neighborhood of $C$ disjoint from $E$, the function
$-m^{-1}\log\|\sigma\|^2$ is therefore pluriharmonic off $C$ and
equals $-\log|w|^2$ plus a smooth function. It is the required
divergent collar.

For condition \textup{(iv)}, consider the restriction
$r:H^1(S,\OO_S)\to H^1(C,\OO_C)$. If $r$ is not injective,
\cite[Proposition~2.7]{CLPT2018} makes a multiple of $C$ a fiber,
and the punctured-base construction applies. Otherwise, since
$q(S)\ge g(C)$, injectivity makes $r$ an isomorphism.
Hodge decomposition and complex
conjugation then give an isomorphism
$H^1(S,\mathbb Q)\to H^1(C,\mathbb Q)$, whose dual is
$j_*\otimes\mathbb Q$. Consequently $K\otimes\mathbb Q=0$.
As a subgroup of the free abelian group $H_1(C,\mathbb Z)$, $K$
is torsion-free, hence $K=0$, and the preceding character argument
applies.
Finally, for \textup{(v)}, \cite[Theorem~1.6]{CFSTZ2026} shows that
an elliptic curve of infinite Ueda type has a unitary flat neighborhood.
Apply
Lemma~\ref{lem:surface-boundary-tools}\textup{(ii)} in all cases.
\end{proof}

\section{Holomorphic convexity for curve fiberings}\label{sec:curve-fibers}
Complex spaces in this section are reduced, Hausdorff, and countable at
infinity. A smooth function is locally the restriction of a smooth
function in an ambient complex manifold. Fiber components refer to
their reduced supports.

\begin{theorem}[=Theorem~\ref{thm:curve-fiber-introduction}]\label{thm:curve-fiber-convexity}
Let $X$ have pure dimension $n\ge2$ and admit a continuous real-valued
psh exhaustion $\rho$. If a holomorphic map
\[
 F:X\longrightarrow\C^{n-1}
 \quad\text{satisfies}\quad
 \dim F^{-1}(z)\le1\qquad(z\in F(X)),
\]
then $X$ is holomorphically convex and admits a smooth psh exhaustion.
\end{theorem}

For smooth $X$, the convexity argument follows
\cite[Corollary~4.2 and pp.~463--464]{VijituForum2022}.
For singular $X$, we first control compact fiber components using
analytic cycles. This allows us to choose exhaustion shells whose
fibers are Stein and then apply the same convexity criteria.

\subsection{Compact fiber components}
Write $F=(f_1,\ldots,f_{n-1})$. On each local irreducible branch of
$X$ at $x$, the fiber through $x$ is defined by the $n-1$ equations
$f_k=f_k(x)$. The analytic dimension theorem therefore gives
\[
 \dim_x F^{-1}(F(x))\ge n-(n-1)=1.
\]
Together with the upper bound in the theorem, this proves that $F$
is equidimensional with one-dimensional fibers on every local branch.
Its graph and projection are
\[
 \iota:X\longrightarrow\Gamma_F\subset\C^{n-1}\times X,
 \qquad x\longmapsto(F(x),x),\qquad
 p:\Gamma_F\longrightarrow\C^{n-1}.
\]
Since $\iota$ is an isomorphism and $p\circ\iota=F$, the reduced graph,
with multiplicity one on each component, defines an analytic family
of geometric fiber cycles over the normal base $\C^{n-1}$
\cite[Chapter~I, \S1(b), p.~816]{BarletMagnusson1998}. Write this family as
\begin{equation}\label{eq:curve-fiber-cycles}
 \mathcal F_z=\sum_jm_j(z)[C_j(z)],\qquad
 m_j(z)\in\mathbb Z_{>0},\qquad
 |\mathcal F_z|=F^{-1}(z).
\end{equation}
The fiber cycles are locally finite, with the zero cycle over an empty fiber.

\begin{lemma}[Uniform area]\label{curve:lem:area}
Fix a Hermitian metric $h$ on $X$. Given compact $L\subset X$ and
$z_0\in\C^{n-1}$, there are a neighborhood $B$ of $z_0$ and $A>0$
such that
\[
 \int_C h\le A
\]
for every compact irreducible component $C\subset L$ of a fiber over $B$.
\end{lemma}
\begin{proof}
Choose $\chi\in C_c^\infty(X)$ with $0\le\chi\le1$ and $\chi=1$
near $L$. The map
\[
 I:\C^{n-1}\longrightarrow\R_{\ge0},\qquad
 I(z)=\int_{\mathcal F_z}\chi h
\]
is continuous by the theorem on integration of compactly supported
forms over analytic families of cycles
\cite[Chapter~IV, Proposition~2.3.1]{BarletMagnusson2014}.
Choose $B$ so that $I(z)<I(z_0)+1$ on $B$ and set $A=I(z_0)+1$.
If $C=C_j(z)\subset L$, its multiplicity in
\eqref{eq:curve-fiber-cycles} is positive. Hence
\[
 \int_C h=\int_C\chi h
 \le\int_{\mathcal F_z}\chi h=I(z)<A.
\]
\end{proof}

\begin{lemma}[Stability]\label{curve:lem:stability}
Let $K\subset X$ be compact. If every irreducible component of
$F^{-1}(z_0)$ meeting $K$ is noncompact, the same holds for all
fibers over a neighborhood of $z_0$.
\end{lemma}
\begin{proof}
Suppose the conclusion fails. There are $z_j\to z_0$, compact
irreducible components $C_j\subset F^{-1}(z_j)$, and
$x_j\in C_j\cap K$. After taking a subsequence, $x_j\to x\in K$.
The maximum principle on the normalization of $C_j$ makes
$\rho|_{C_j}$ constant. Thus, with
$R=\max_K\rho$,
\[
 \rho|_{C_j}=\rho(x_j)\le R,\qquad
 C_j\subset L:=\{\rho\le R\}\Subset X.
\]
Lemma~\ref{curve:lem:area} bounds $\int_{C_j}h$ for all sufficiently
large $j$. Bishop compactness for cycles on complex spaces
\cite[Theorem~1]{Barlet1999} yields a convergent subsequence
\[
 [C_j]\longrightarrow T\quad\text{in the space of compact
 effective one-cycles of }X.
\]
Convergence of compact cycles gives upper semicontinuity of their
supports, so $x_j\to x$ implies $x\in|T|$. Convergence of the associated
currents gives $|T|\subset L$. Moreover, a point outside
$F^{-1}(z_0)$ has a neighborhood disjoint from $F^{-1}(z_j)$ for all
large $j$, so it is outside $|T|$. Therefore
\[
 x\in|T|\subset L\cap F^{-1}(z_0).
\]
An irreducible component of $|T|$ through $x$ is a compact curve
in the pure one-dimensional fiber $F^{-1}(z_0)$ and hence a component
of that fiber. It meets $K$, contradicting the hypothesis.
\end{proof}

\subsection{Convexity and reduction}
We use the following three criteria for complex spaces
\cite[Propositions~2.2, 2.4 and Corollary~2.5]{VijituForum2022}.
\begin{enumerate}[label=\textup{(\roman*)},leftmargin=2.5em]
\item A complex space with a continuous psh function $u$ is
holomorphically convex if and only if every $\{u<t\}$ is.
\item A space with a continuous psh exhaustion is holomorphically convex
if it is locally holomorphically convex over a Stein base.
\item If $\pi:M\to S$ is holomorphic, $S$ is Stein, and
an open set $D\subset M$ admits a continuous psh exhaustion, then
$D$ is holomorphically convex provided $\overline D^{\,M}\to S$ is
proper and a neighborhood
of $\partial_M D$ has Stein reduced fibers.
\end{enumerate}
Here local holomorphic convexity over $S$ means that
every point of $S$ has a neighborhood $V$ whose inverse image is
holomorphically convex.

For \textup{(iii)}, the Stein-neighborhood construction in the cited
proof gives, over a relatively compact base neighborhood $V\Subset S$,
finitely many smooth functions $h_j$
and nonnegative base cutoffs $\tau_j$ for which
\[
 h_0=\sum_j(\tau_j\circ\pi)h_j,\qquad
 h=h_0+\Lambda\phi\circ\pi
\]
is strictly psh on a collar $W$ of the boundary over $\overline V$,
where $\phi$ is strictly psh on the base. The product rule and
Cauchy--Schwarz on the finitely many local embeddings give constants $a,b,c>0$ such that
\[
 \Levi_{h_0}(\xi)\ge a\|\xi\|^2-b\|d\pi(\xi)\|^2,
 \qquad
 \Levi_{\phi\circ\pi}(\xi)\ge c\|d\pi(\xi)\|^2.
\]
Taking $\Lambda c>b$ makes $h$ strictly psh on the collar.
The functions are smooth on a neighborhood of
$\overline D\cap\pi^{-1}(\overline V)$. Let $\eta_D$ be a continuous
psh exhaustion of $D$ and put
\[
 \psi=\eta_D+h\quad\text{on }D\cap\pi^{-1}(V).
\]
The set
\[
 H=\bigl(\overline D\cap\pi^{-1}(\overline V)\bigr)\setminus W
\]
is a compact subset of $D$. Choose a constant $A$ larger than
$\max_H\psi$ when $H$ is nonempty. Then $\{\psi>A\}\subset W$,
where $\psi$ is strongly psh. For every compact $T\subset V$, the
function $h$ has a lower bound $m_T$ on
$\overline D\cap\pi^{-1}(T)$, and
\[
 \{\psi\le t\}\cap\pi^{-1}(T)
 \subset\{\eta_D\le t-m_T\}\cap\pi^{-1}(T)\Subset D.
\]
The set on the left is closed in the compact set on the right.
Thus $\psi$ is a vertical exhaustion and defines a $1$-convex map.
Holomorphic convexity follows from
\cite[Theorem~2.3 and Proposition~2.4]{VijituForum2022}.

\begin{lemma}\label{curve:lem:neighborhood}
For every $\lambda\in\R$ and $z_0\in\C^{n-1}$, there are $c>\lambda$
and a ball $B=B_r(z_0)$ such that
\[
 D=F^{-1}(B)\cap\{\rho<c\}
 \quad\text{and}\quad
 F^{-1}(B)\cap\{\rho<\lambda\}
\]
are holomorphically convex.
\end{lemma}
\begin{proof}
Let $\mathcal V_0$ be the set of constant $\rho$-values on compact irreducible
components of $F^{-1}(z_0)$. For each $R$, every component whose
value is at most $R$ lies in the compact set $\{\rho\le R\}$.
The locally finite family of irreducible components of the central
fiber has only finitely many members meeting this set. Therefore
\[
 \#\bigl(\mathcal V_0\cap(-\infty,R]\bigr)<\infty.
\]
Choose $\lambda<a<c<b$ with $[a,b]\cap\mathcal V_0=\varnothing$ and put
$K=\{a\le\rho\le b\}$. This is compact, and no compact component
of $F^{-1}(z_0)$ meets it. Lemma~\ref{curve:lem:stability} gives a
ball $B=B_r(z_0)$ such that no compact component of a fiber over
$B$ meets $K$. Set
\[
 M=F^{-1}(B),\qquad D=M\cap\{\rho<c\},\qquad
 U=M\cap\{a<\rho<b\}.
\]
Continuity of $\rho$ gives
$\partial_M D\subset\{\rho=c\}\cap M\subset U$.
If a reduced fiber of $F|_U$ had a compact irreducible component $C$,
compactness would make $C$ a closed analytic subset of $X$. Since
it has dimension one, it would be a component of the full $F$-fiber.
This contradicts the choice of $B$, because $C\subset K$.
Hence every reduced $U$-fiber is Stein by the criterion for curves
without compact irreducible components
\cite[\S3, p.~1663]{Vajaitu2022}.

On $D$, define
\[
 \eta=\max\left\{\frac1{c-\rho},
                  \frac1{r^2-\|F-z_0\|^2}\right\}.
\]
The function $t\mapsto(\alpha-t)^{-1}$ is convex and increasing on
$t<\alpha$, so both terms and their maximum are continuous and psh.
For $t>0$,
\[
 \{\eta\le t\}
 =\{\rho\le c-1/t\}\cap
 F^{-1}\{\|z-z_0\|^2\le r^2-1/t\}\Subset D.
\]
Thus $\eta$ is a continuous psh exhaustion of $D$. For compact $T\subset B$,
\[
 \overline D^{\,M}\cap F^{-1}(T)
 \subset\{\rho\le c\}\cap F^{-1}(T).
\]
The left-hand side is closed in the compact set on the right, so
$\overline D^{\,M}\to B$ is proper. Criterion~\textup{(iii)} now
gives holomorphic convexity of $D$. Since $\lambda<c$,
\[
 F^{-1}(B)\cap\{\rho<\lambda\}
 =\{x\in D:\rho(x)<\lambda\},
\]
and criterion~\textup{(i)} gives the second assertion.
\end{proof}

\begin{proof}[Proof of Theorem~\ref{thm:curve-fiber-convexity}]
For $X(\lambda)=\{\rho<\lambda\}$,
Lemma~\ref{curve:lem:neighborhood} makes
$F|_{X(\lambda)}:X(\lambda)\to\C^{n-1}$ locally holomorphically
convex. The continuous psh function
$\eta_\lambda=(\lambda-\rho)^{-1}$ exhausts $X(\lambda)$, since
\[
 \{\eta_\lambda\le t\}=\{\rho\le\lambda-1/t\}
 \Subset X(\lambda)\qquad(t>0).
\]
Criterion~\textup{(ii)} makes every $X(\lambda)$ holomorphically
convex, and criterion~\textup{(i)} then applies to $\rho$ on $X$.
Let $q:X\to Z$ be the proper Remmert reduction onto its Stein
space $Z$, and let $u$ be a smooth strictly psh exhaustion of $Z$.
Holomorphic pullback preserves smoothness in local embeddings and
plurisubharmonicity, while properness gives
\[
 u\circ q\in C^\infty(X)\cap\operatorname{PSH}(X),\qquad
 \{u\circ q\le t\}=q^{-1}(\{u\le t\})\Subset X.\qedhere
\]
\end{proof}

\begin{corollary}[Simultaneous contraction]\label{cor:curve-contraction}
The Remmert reduction $q:X\to Z$ has connected fibers,
$q_*\OO_X=\OO_Z$, and $F=\widetilde F\circ q$ for a holomorphic
$\widetilde F:Z\to\C^{n-1}$. Its nontrivial fibers are exactly the
maximal connected unions of compact irreducible fiber curves of $F$.
\end{corollary}
\begin{proof}
By $q_*\OO_X=\OO_Z$, each coordinate $f_k$ of $F$ descends to a
holomorphic function $\widetilde f_k$ on $Z$. These functions define
$\widetilde F$ with $F=\widetilde F\circ q$.
Every $q$-fiber is therefore a compact connected analytic subset of
an $F$-fiber. If it is not a point, the dimension bound and local
finiteness make it a finite connected union of compact irreducible
components of that $F$-fiber.

Conversely, the image of a compact irreducible curve in the Stein
space $Z$ is a point, so $q$ contracts each such curve. Consider a connected union $E$ of
compact irreducible components of one $F$-fiber. The normalization
maximum principle makes $\rho$ constant on each component.
Intersecting components have the same value. By local finiteness and
connectedness, any two components can be joined by a finite chain of
intersecting components.
Thus $\rho|_E=a$ for one $a\in\R$, so $E\subset\{\rho\le a\}$.
Only finitely many fiber components meet this compact set. Hence
$E$ is a finite connected union and $q$ is constant on $E$.
It follows that the maximal such unions are precisely the nontrivial
$q$-fibers.
\end{proof}

% Alphabetic bibliography, formatted with BibTeX alpha.bst.
\newcommand{\etalchar}[1]{$^{#1}$}
\begin{thebibliography}{CLPT18}

\bibitem[Bar99]{Barlet1999}
D.~Barlet.
\newblock How to use the cycle space in complex geometry.
\newblock In M.~Schneider and Y.-T. Siu, editors, {\em Several Complex Variables}, volume~37 of {\em MSRI Publications}, pages 25--42. Cambridge University Press, 1999.
\newblock \url{https://library.slmath.org/books/Book37/files/barlet.pdf}.

\bibitem[Bar15]{Barlet2015}
D.~Barlet.
\newblock Some examples due to {H.~Hironaka}, 2015.
\newblock Preprint, \href{https://arxiv.org/abs/1501.01129}{arXiv:1501.01129v2}, Section~6.

\bibitem[BM98]{BarletMagnusson1998}
D.~Barlet and J.~Magn\'usson.
\newblock {Int\'egration de classes de cohomologie m\'eromorphes et diviseurs d'incidence}.
\newblock {\em Ann. Sci. \'Ecole Norm. Sup. (4)}, 31(6):811--842, 1998.
\newblock \href{https://doi.org/10.1016/S0012-9593(99)80004-4}{doi:\nolinkurl{10.1016/S0012-9593(99)80004-4}}. \url{https://www.numdam.org/item/ASENS_1998_4_31_6_811_0.pdf}.

\bibitem[BM14]{BarletMagnusson2014}
D.~Barlet and J.~Magn\'usson.
\newblock {\em {Cycles analytiques complexes I: th\'eor\`emes de pr\'eparation des cycles}}, volume~22 of {\em Cours Sp\'ecialis\'es}.
\newblock Soci\'et\'e Math\'ematique de France, Paris, 2014.

\bibitem[CFS{\etalchar{+}}26]{CFSTZ2026}
Y.~Chen, S.~Filip, S.~Sun, V.~Tosatti, and J.~Zhang.
\newblock Compact {K\"ahler} surfaces with semipositive anticanonical bundle, 2026.
\newblock Preprint, \href{https://arxiv.org/abs/2609.26716v1}{arXiv:2609.26716v1}, submitted September~22, 2026.

\bibitem[CLPT18]{CLPT2018}
B.~Claudon, F.~Loray, J.~V. Pereira, and F.~Touzet.
\newblock Compact leaves of codimension one holomorphic foliations on projective manifolds.
\newblock {\em Ann. Sci. \'Ecole Norm. Sup. (4)}, 51(6):1457--1506, 2018.
\newblock \href{https://doi.org/10.24033/asens.2379}{doi:\nolinkurl{10.24033/asens.2379}}. \href{https://arxiv.org/abs/1512.06623}{Author preprint, Theorems~1.3, 2.5, 4.4 and Proposition~2.7}.

\bibitem[Dem82a]{DemaillyL2}
J.-P. Demailly.
\newblock Estimations $L^2$ pour l'op\'erateur $\barp$ d'un fibr\'e
vectoriel holomorphe semi-positif au-dessus d'une vari\'et\'e
k\"ahl\'erienne compl\`ete.
\newblock \emph{Ann. Sci. \'Ecole Norm. Sup. (4)}, 15(3):457--511, 1982.
\newblock \href{https://doi.org/10.24033/asens.1434}{doi:\nolinkurl{10.24033/asens.1434}}.

\bibitem[Dem82b]{Demailly1982}
J.-P. Demailly.
\newblock \emph{Sur diff\'erents aspects de la positivit\'e en analyse complexe}.
\newblock Th\`ese de doctorat d'\'Etat, Universit\'e Pierre et Marie Curie
(Paris~VI), 1982.
\newblock \href{https://www-fourier.ujf-grenoble.fr/~demailly/manuscripts/these_etat_intro.pdf}{Thesis introduction}.

\bibitem[DNS22]{DNS22}
T.-C. Dinh, V.-A. Nguy\^en, and N.~Sibony.
\newblock Unique ergodicity for foliations on compact {K\"ahler} surfaces.
\newblock {\em Duke Math. J.}, 171(13):2627--2698, 2022.
\newblock \href{https://doi.org/10.1215/00127094-2022-0044}{doi:\nolinkurl{10.1215/00127094-2022-0044}}. \href{https://arxiv.org/abs/1811.07450v2}{Author version with the cited equation numbers and Appendix~B}.

\bibitem[Eno81]{Enoki}
I.~Enoki.
\newblock Surfaces of class {$\mathrm{VII}_0$} with curves.
\newblock {\em Tohoku Math. J. (2)}, 33(4):453--492, 1981.
\newblock \href{https://doi.org/10.2748/tmj/1178229349}{doi:\nolinkurl{10.2748/tmj/1178229349}}. \href{https://www.jstage.jst.go.jp/article/tmj1949/33/4/33_4_453/_pdf/-char/en}{Original article, pp.~453--454}.

\bibitem[FT18]{FilipTosatti2018}
S.~Filip and V.~Tosatti.
\newblock Smooth and rough positive currents.
\newblock {\em Annales de l'Institut Fourier}, 68(7):2981--2999, 2018.
\newblock \href{https://doi.org/10.5802/aif.3234}{doi:\nolinkurl{10.5802/aif.3234}}.

\bibitem[FN05]{FN2005}
Y.~Fujimoto and N.~Nakayama.
\newblock Compact complex surfaces admitting non-trivial surjective endomorphisms.
\newblock {\em Tohoku Math. J. (2)}, 57(3):395--426, 2005.
\newblock \href{https://doi.org/10.2748/tmj/1128703004}{doi:\nolinkurl{10.2748/tmj/1128703004}}. \href{https://www.kurims.kyoto-u.ac.jp/preprint/file/RIMS1431.pdf}{Author preprint, Section~2 and Lemma~2.1}.

\bibitem[Gra62]{Grauert1962}
Hans Grauert.
\newblock {{\"U}ber Modifikationen und exzeptionelle analytische Mengen}.
\newblock {\em Mathematische Annalen}, 146(4):331--368, 1962.
\newblock \href{https://doi.org/10.1007/BF01441136}{doi:\nolinkurl{10.1007/BF01441136}}.

\bibitem[Koi21]{Koike2021}
Takayuki Koike.
\newblock Linearization of transition functions of a semi-positive line bundle along a certain submanifold.
\newblock {\em Annales de l'Institut Fourier}, 71(5):2237--2271, 2021.
\newblock \href{https://doi.org/10.5802/aif.3439}{doi:\nolinkurl{10.5802/aif.3439}}.

\bibitem[KU25]{KoikeUeda2025}
Takayuki Koike and Tetsuo Ueda.
\newblock Plurisubharmonic functions on affine line bundles over compact {K\"ahler} manifolds.
\newblock {\em Publ. Res. Inst. Math. Sci.}, 61(1):139--152, 2025.
\newblock \href{https://doi.org/10.4171/PRIMS/61-1-3}{doi:\nolinkurl{10.4171/PRIMS/61-1-3}}.

\bibitem[MS20]{MongodiSlodkowski2020}
S.~Mongodi and Z.~S{\l}odkowski.
\newblock Domains with a continuous exhaustion in weakly complete surfaces.
\newblock {\em Math. Z.}, 296(3--4):1011--1019, 2020.
\newblock \href{https://doi.org/10.1007/s00209-020-02466-z}{doi:\nolinkurl{10.1007/s00209-020-02466-z}}. \href{https://arxiv.org/abs/1904.03497}{arXiv:1904.03497}.

\bibitem[Ohs81]{Ohsawa1981}
T.~Ohsawa.
\newblock Weakly 1-complete manifold and {Levi} problem.
\newblock {\em Publ. Res. Inst. Math. Sci.}, 17(1):153--164, 1981.
\newblock \href{https://doi.org/10.2977/PRIMS/1195186709}{doi:\nolinkurl{10.2977/PRIMS/1195186709}}.

\bibitem[Ohs21]{Ohsawa2021}
T.~Ohsawa.
\newblock {Nishino's} rigidity, locally pseudoconvex maps, and holomorphic motions.
\newblock {\em RIMS K\^oky\^uroku}, 2175:27--46, 2021.
\newblock \href{https://hdl.handle.net/2433/263965}{Kyoto University repository}.
\href{https://www.kurims.kyoto-u.ac.jp/~kyodo/kokyuroku/contents/pdf/2175-03.pdf}{Author text, Section~6}.

\bibitem[Sta]{StacksHironaka}
{The Stacks Project Authors}.
\newblock The {Stacks} project.
\newblock Lemma~110.67.1, Tag~08J1, ``Being projective is not local on the base.'' \url{https://stacks.math.columbia.edu/tag/08J1}.

\bibitem[Ued83]{Ueda1983}
Tetsuo Ueda.
\newblock On the neighborhood of a compact complex curve with topologically trivial normal bundle.
\newblock {\em Journal of Mathematics of Kyoto University}, 22(4):583--607, 1982/83.
\newblock \href{https://doi.org/10.1215/kjm/1250521670}{doi:\nolinkurl{10.1215/kjm/1250521670}}.

\bibitem[V{\^\i}j22a]{Vajaitu2022}
V.~V{\^\i}j{\^\i}itu.
\newblock Holomorphic convexity and pseudoconvex surfaces.
\newblock {\em Ann. Sc. Norm. Super. Pisa Cl. Sci. (5)}, 23(4):1659--1678, 2022.
\newblock \href{https://doi.org/10.2422/2036-2145.202011_018}{doi:\nolinkurl{10.2422/2036-2145.202011_018}}. \url{https://journals.sns.it/index.php/annaliscienze/article/view/5700}.

\bibitem[V{\^\i}j22b]{VijituForum2022}
V.~V{\^\i}j{\^\i}itu.
\newblock Holomorphic convexity of pseudoconvex spaces in terms of the rank of structural sheaf.
\newblock {\em Forum Math.}, 34(2):455--468, 2022.
\newblock \href{https://doi.org/10.1515/forum-2021-0262}{doi:\nolinkurl{10.1515/forum-2021-0262}}.

\bibitem[Wak80]{Wakabayashi1980}
I.~Wakabayashi, editor.
\newblock Collected problems around several complex variables (in Japanese).
\newblock {\em Sugaku}, 32(2):161--187, 1980.
\newblock \href{https://doi.org/10.11429/sugaku1947.32.161}{doi:\nolinkurl{10.11429/sugaku1947.32.161}}.
\newblock S.~Nakano's contribution: Section~I.1, Problem~1, p.~162.

\bibitem[W{\l}o05]{Wlodarczyk05}
J.~W{\l}odarczyk.
\newblock Simple {Hironaka} resolution in characteristic zero.
\newblock {\em J. Amer. Math. Soc.}, 18(4):779--822, 2005.
\newblock \href{https://doi.org/10.1090/S0894-0347-05-00493-5}{doi:\nolinkurl{10.1090/S0894-0347-05-00493-5}}. Author version: \href{https://arxiv.org/abs/math/0401401}{arXiv:math/0401401}.

\bibitem[Wu26]{Wu2026}
X.~Wu.
\newblock Generalized {Ueda} obstruction classes and applications to non-semi-positivity of line bundles, 2026.
\newblock Preprint, \href{https://arxiv.org/abs/2606.22821v1}{arXiv:2606.22821v1}.

\bibitem[Zha25]{Zhang2025}
Yangyang Zhang.
\newblock On the non-semipositivity of a nef and big line bundle on {Grauert's} example, 2025.
\newblock \href{https://arxiv.org/abs/2512.23446v1}{arXiv:2512.23446v1}, submitted 29 December 2025.

\end{thebibliography}
\end{document}
